\chapter{Proper Morphisms, Ambidexterity, and Compactification}
\label{chap:proper-ambidexterity-compactification}

\section{Purpose, geometric scope, and standing conventions}
\label{sec:proper-purpose}

Chapter~\ref{chap:stable-brave-new-motivic} constructed, for every
quasi-compact quasi-separated spectral scheme \(S\), a compactly generated
stable presentably symmetric monoidal category
\[
    \SH(S),
\]
together with adjunctions
\[
    f^*:\SH(S)\rightleftarrows\SH(X):f_*
\]
for arbitrary morphisms
\[
    f:X\longrightarrow S
\]
between quasi-compact quasi-separated spectral schemes, and adjunctions
\[
    p_\sharp:\SH(X)\rightleftarrows\SH(S):p^*
\]
for smooth morphisms \(p\) of finite presentation.  It also supplied stable
localization, closed recollement, smooth and closed exchange, stable homotopy
purity, and virtual Thom twists
\[
    \Sigma^\xi,
    \qquad
    \xi\in K(X).
\]

The purpose of the present chapter is to establish the proper geometric input
required for exceptional functoriality.  No new direct-image functor is needed
for a proper morphism: the functor \(f_*\) already exists as the right adjoint
of inverse image.  The new assertions are that arbitrary direct image preserves small
colimits, proper direct image satisfies arbitrary base change and the
projection formula, proper direct image satisfies open--proper exchange and
the resulting support property, and smooth proper direct image is identified
with smooth direct image after the tangent Thom twist.  The chapter then
proves the compactification theorem and the contractibility of the category
of compactifications.  These results are the precise geometric hypotheses
used in the construction of \(f_!\) and \(f^!\) in the next chapter.

\begin{convention}[Geometric classes]
\label{conv:proper-geometric-classes}
Throughout this chapter, all spectral schemes are quasi-compact and
quasi-separated.  We use the following classes of morphisms.

\begin{enumerate}
    \item \(\mathcal E\) is the class of separated morphisms of finite
    presentation.

    \item \(\mathcal P\) is the class of morphisms
    \[
        p:Y\longrightarrow X
    \]
    whose classical truncation
    \[
        p_{\mathrm{cl}}:
        Y_{\mathrm{cl}}
        \longrightarrow
        X_{\mathrm{cl}}
    \]
    is proper.

    \item \(\mathcal J\subseteq\mathcal E\) is the class of quasi-compact open
    immersions.

    \item \(\mathcal S\subseteq\mathcal E\) is the class of smooth separated
    morphisms of finite presentation.
\end{enumerate}

A compactification always means a factorization by a morphism in
\(\mathcal J\) followed by a morphism in \(\mathcal P\).  No spectral
finite-presentation condition is imposed on the proper part of a
compactification.
\end{convention}

\begin{convention}[Properness and projective presentations]
\label{conv:spectral-properness}
A morphism
\[
    p:Y\longrightarrow X
\]
belongs to \(\mathcal P\) when its classical truncation
\[
    p_{\mathrm{cl}}:
    Y_{\mathrm{cl}}
    \longrightarrow
    X_{\mathrm{cl}}
\]
is proper.

When \(p\) is of finite presentation, this condition is equivalent to
properness in the spectral-geometric sense.  Thus, for a morphism of finite
presentation,
\[
    p\text{ is spectrally proper}
    \quad\Longleftrightarrow\quad
    p_{\mathrm{cl}}\text{ is proper}.
\]
Separatedness of a morphism in \(\mathcal E\) is likewise detected on
classical truncations.

Projectivity is defined below by the existence of a finite spectral
projective presentation whose individual stages are actual closed embeddings
into projective bundles associated to finite locally free modules.  Every
projective morphism is of finite presentation and belongs to \(\mathcal P\).
Classical truncation carries every spectral projective presentation to a
classical projective presentation.  No converse projectivity-detection
statement is used in this chapter.
\end{convention}

\begin{convention}[Tangent and normal Thom classes]
\label{conv:tangent-normal-classes}
Let
\[
    i:Z\hookrightarrow X
\]
be a quasi-smooth closed immersion whose conormal module is finite locally
free.  Its normal Thom class is
\[
    \nu_i
    :=
    [N^*_{Z/X}]
    \in
    K(Z).
\]
Under the coordinate-module convention of
Remark~\ref{rem:thom-dual-convention}, the twist \(\Sigma^{\nu_i}\) is the
Thom suspension by the geometric normal bundle
\[
    \mathbf N_i
    =
    \mathbb V_Z(N^*_{Z/X}).
\]

Let
\[
    p:X\longrightarrow S
\]
be smooth of finite presentation.  Its diagonal is a quasi-smooth
immersion; when \(p\) is separated, it is a closed immersion.  In either
case its conormal module is
\[
    N^*_{\Delta_p}
    \simeq
    \Omega_{X/S}.
\]
We define the tangent Thom class of \(p\) by
\[
    T_p
    :=
    [\Omega_{X/S}]
    \in
    K(X).
\]
Thus \(\Sigma^{T_p}\) is the Thom suspension by the geometric relative
tangent bundle.  This convention incorporates the dualization already built
into the definition
\[
    \mathbb V_X(E)
    =
    \operatorname{Spec}_X\operatorname{Sym}(E).
\]
\end{convention}

\begin{remark}[Use of ordinary comparison]
\label{rem:proper-comparison-method}
Chapter~\ref{chap:stable-brave-new-motivic} constructed natural symmetric
monoidal equivalences
\[
    \Theta_S^{\mathrm{st}}:
    \SH(S)
    \xrightarrow{\simeq}
    \SH_{\mathrm{cl}}(S_{\mathrm{cl}})
\]
compatible with inverse image, arbitrary direct image, smooth direct image,
closed direct image, exceptional inverse image for closed immersions,
localization, homotopy purity, and Thom twists.

The Beck--Chevalley, projection, support, and ambidexterity transformations
used below are constructed internally from the stable coefficient system and
its adjunctions.  Their invertibility is sometimes detected after applying
stable classical comparison.

The coefficientwise relative closed-purity transformation is treated
slightly differently.  The preceding chapter supplied stable homotopy purity
for the associated Thom object, but did not construct coefficientwise
relative purity on arbitrary objects.  We therefore characterize the
relative closed-purity transformation below by transport of the canonical
ordinary transformation through
\(\Theta^{\mathrm{st}}\).  This uses only the already-constructed functor
\(i^!\) for a closed immersion.  It does not define any exceptional operation
for a general separated morphism.
\end{remark}

\subsection{The chapter-level output}

The constructions below culminate in the following assertions.
\begin{enumerate}
    \item Spectral projective spaces and projective bundles have their usual
    smoothness, properness, base-change, and classical-truncation properties.

    \item For every smooth proper morphism \(p:X\to S\), there is a canonical
    ambidexterity equivalence
    \[
        p_\sharp\Sigma^{-T_p}
        \xrightarrow{\simeq}
        p_*.
    \]

    \item For every morphism \(f\), arbitrary direct image
    \[
        f_*:\SH(X)\longrightarrow\SH(S)
    \]
    preserves small colimits.

    \item Proper direct image satisfies arbitrary base change and the proper
    projection formula.

    \item Proper morphisms satisfy open--proper exchange.  In particular,
    they satisfy the support property for proper modifications that are
    equivalences over an open complement.

    \item Every separated morphism of finite presentation admits a spectral
    compactification.  Its \(\infty\)-category of compactifications is
    nonempty and cofiltered, and hence has contractible groupoid completion.

    \item The compactification diagram
    \[
        c=(X\xrightarrow{j}\overline X\xrightarrow{p}S)
        \longmapsto
        p_*j_\sharp
    \]
    takes values in colimit-preserving functors and carries every refinement
    to an equivalence through coherently compatible comparison maps.

    \item For every finite composable string of morphisms in \(\mathcal E\),
    the \(\infty\)-category of compactification grids is nonempty and
    cofiltered, and hence has contractible groupoid completion.  The
    associated open--proper exchange equivalences identify all iterated
    compactification formulas coherently.
\end{enumerate}

\section{Spectral projective geometry}
\label{sec:spectral-projective-geometry}

\subsection{Spectral projective space}

\begin{construction}[Spectral projective space]
\label{constr:spectral-projective-space}
Let \(S\) be a spectral scheme and let \(n\geq0\).  For
\(0\leq a\leq n\), let
\[
    U_a
    \simeq
    \mathbb A^n_S
\]
be a copy of brave new affine \(n\)-space.  For \(a\neq b\), identify the
principal open in \(U_a\) on which the coordinate corresponding to \(b\) is
invertible with the corresponding principal open in \(U_b\) by the usual
homogeneous-coordinate formulas.  These identifications satisfy the cocycle
condition on triple intersections.  Spectral Zariski gluing therefore
produces a spectral scheme
\[
    \mathbb P^n_S
\]
with an open cover by the \(U_a\).
\end{construction}

\begin{proposition}[Base change for projective space]
\label{prop:projective-space-base-change}
For every morphism of spectral schemes
\[
    f:T\longrightarrow S,
\]
there is a canonical equivalence
\[
    \mathbb P^n_T
    \simeq
    \mathbb P^n_S\times_ST.
\]
The equivalence is coherent in \(f\) and preserves the standard affine
charts and their homogeneous coordinates.
\end{proposition}

\begin{proof}
Base change carries brave new affine space to brave new affine space and
preserves principal open immersions and their intersections.  Therefore the
pullback of the standard affine atlas of
\(\mathbb P^n_S\) is canonically identified with the standard affine atlas
used to construct \(\mathbb P^n_T\), including all transition
isomorphisms and cocycle data.

Both
\[
    \mathbb P^n_S\times_ST
\]
and
\[
    \mathbb P^n_T
\]
are consequently obtained by gluing the same spectral affine charts along
the same principal-open identifications.  The uniqueness clause in spectral
Zariski descent gives a canonical equivalence
\[
    \mathbb P^n_T
    \simeq
    \mathbb P^n_S\times_ST.
\]
Coherence in \(f\) follows from the coherence of base change for affine
spaces, principal opens, and their descent data.
\end{proof}

\begin{proposition}[Classical truncation of projective space]
\label{prop:projective-space-truncation}
There is a canonical equivalence of ordinary schemes
\[
    (\mathbb P^n_S)_{\mathrm{cl}}
    \simeq
    \mathbb P^n_{S_{\mathrm{cl}}}.
\]
It identifies the standard affine covers and is natural in \(S\).
\end{proposition}

\begin{proof}
Classical truncation carries
\(\mathbb A^n_S\) to \(\mathbb A^n_{S_{\mathrm{cl}}}\), preserves the
principal opens used in the gluing construction, and preserves their
Cartesian intersections.  Truncating the standard spectral gluing diagram
therefore gives the standard classical gluing diagram for projective space.
The asserted equivalence follows from ordinary Zariski descent.
\end{proof}

\begin{proposition}[Smoothness and properness of projective space]
\label{prop:projective-space-smooth-proper}
The structural morphism
\[
    \pi:\mathbb P^n_S\longrightarrow S
\]
is smooth, proper, and of finite presentation.
\end{proposition}

\begin{proof}
Smoothness and finite presentation are local on the source.  On each standard
chart, \(\pi\) is the projection
\[
    \mathbb A^n_S\longrightarrow S,
\]
which is smooth and of finite presentation.  Hence \(\pi\) is smooth and of
finite presentation globally.

By Proposition~\ref{prop:projective-space-truncation}, the classical
truncation of \(\pi\) is the usual projective morphism
\[
    \mathbb P^n_{S_{\mathrm{cl}}}
    \longrightarrow
    S_{\mathrm{cl}},
\]
which is proper.  Properness is detected on classical truncations under the
finite-presentation hypothesis, so \(\pi\) is proper.
\end{proof}

\begin{construction}[The twisting line bundle]
\label{constr:spectral-O-one}
The standard transition functions on the pairwise intersections of the
affine charts glue the trivial line bundles to an invertible
\(\mathcal O_{\mathbb P^n_S}\)-module
\[
    \mathcal O_{\mathbb P^n_S}(1).
\]
Its dual is denoted
\[
    \mathcal O_{\mathbb P^n_S}(-1).
\]

The construction is compatible with arbitrary base change.  More precisely,
for every morphism
\[
    f:T\longrightarrow S,
\]
the canonical base-change equivalence
\[
    \mathbb P^n_T
    \xrightarrow{\simeq}
    \mathbb P^n_S\times_ST
\]
identifies
\[
    \mathcal O_{\mathbb P^n_T}(1)
    \simeq
    \operatorname{pr}_1^*
    \mathcal O_{\mathbb P^n_S}(1).
\]

Using
Proposition~\ref{prop:projective-space-truncation},
let
\[
    \iota_{\mathbb P^n_S}:
    \mathbb P^n_{S_{\mathrm{cl}}}
    \simeq
    (\mathbb P^n_S)_{\mathrm{cl}}
    \longrightarrow
    \mathbb P^n_S
\]
denote the classical truncation immersion.  Then there is a canonical
equivalence
\[
    \iota_{\mathbb P^n_S}^*
    \mathcal O_{\mathbb P^n_S}(1)
    \simeq
    \mathcal O_{\mathbb P^n_{S_{\mathrm{cl}}}}(1).
\]
\end{construction}

\subsection{The projective-line model of the Thom sphere}

\begin{proposition}[The spectral projective line and the Thom sphere]
\label{prop:spectral-P1-T}
Let
\[
    \infty_S:S\longrightarrow\mathbb P^1_S
\]
be the point at infinity.  There is a canonical equivalence
\[
    (\mathbb P^1_S,\infty_S)
    \simeq
    T_S
\]
in \(\HH_\bullet(S)\).  It is natural under arbitrary base change in
\(S\).
\end{proposition}

\begin{proof}
Write
\[
    \mathbb P^1_S
    =
    U\cup V,
\]
where
\[
    U=\mathbb A^1_0,
    \qquad
    V=\mathbb A^1_\infty,
    \qquad
    U\cap V=\Gm{}_S,
\]
and where the point at infinity belongs to \(V\).

There is a cofiber sequence in the pointed motivic category
\[
    (V,\infty_S)
    \longrightarrow
    (\mathbb P^1_S,\infty_S)
    \longrightarrow
    \mathbb P^1_S/V.
\]
The pointed motivic space
\[
    (V,\infty_S)
    \simeq
    (\mathbb A^1_S,0_S)
\]
is contractible by affine-line invariance.  Hence the quotient map induces a
canonical equivalence
\[
    (\mathbb P^1_S,\infty_S)
    \xrightarrow{\simeq}
    \mathbb P^1_S/V.
\]

Zariski descent for the open cover
\[
    \mathbb P^1_S=U\cup V
\]
gives a pushout square
\[
\begin{tikzcd}[column sep=large,row sep=large]
    \Mot_S(U\cap V) \arrow[r] \arrow[d]
        & \Mot_S(U) \arrow[d] \\
    \Mot_S(V) \arrow[r]
        & \Mot_S(\mathbb P^1_S).
\end{tikzcd}
\]
Taking the cofiber of the lower horizontal morphism identifies
\[
    \mathbb P^1_S/V
    \simeq
    U/(U\cap V).
\]
Using
\[
    U\simeq\mathbb A^1_S
    \qquad\text{and}\qquad
    U\cap V\simeq\Gm{}_S,
\]
we obtain
\[
    \mathbb P^1_S/V
    \simeq
    \mathbb A^1_S/\Gm{}_S
    =
    T_S.
\]
Combining the two equivalences gives
\[
    (\mathbb P^1_S,\infty_S)
    \simeq
    T_S.
\]

The standard open cover, its intersection, the point at infinity, the
Zariski pushout, and affine-line contraction are all preserved by arbitrary
base change.  The equivalence is therefore natural in \(S\).
\end{proof}

\begin{corollary}[Projective-line stabilization]
\label{cor:spectral-P1-stabilization}
There is a canonical symmetric monoidal equivalence
\[
    \SH(S)
    \simeq
    \HH_\bullet(S)
    \bigl[
      (\mathbb P^1_S,\infty_S)^{-1}
    \bigr].
\]
\end{corollary}

\begin{proof}
Apply symmetric monoidal periodization to the equivalence of
Proposition~\ref{prop:spectral-P1-T}.
\end{proof}

\subsection{Projective bundles}

\begin{definition}[Line quotient]
\label{def:spectral-line-quotient}
Let \(T\) be a spectral scheme and let \(M\) be a connective
quasi-coherent \(\mathcal O_T\)-module.

A line quotient of \(M\) is a pair
\[
    (L,q)
\]
consisting of a finite locally free quasi-coherent
\(\mathcal O_T\)-module \(L\) of rank one and a morphism
\[
    q:M\longrightarrow L
\]
such that
\[
    \pi_0(q):
    \pi_0(M)
    \longrightarrow
    \pi_0(L)
\]
is surjective.

Equivalently, Zariski-locally on \(T\), the module \(L\) is trivial and the
morphism \(q\) admits an \(\mathcal O_T\)-linear section.
\end{definition}

\begin{convention}[Rank-zero projective bundles]
\label{conv:rank-zero-projective-bundle}
For every spectral scheme \(T\), put
\[
    \mathbb P^{-1}_T
    :=
    \varnothing.
\]
The unique quasi-coherent module on the empty spectral scheme is denoted
\(\mathcal O_{\mathbb P^{-1}_T}(1)\), and the unique morphism
\[
    \mathcal O_{\mathbb P^{-1}_T}^{\oplus0}
    \longrightarrow
    \mathcal O_{\mathbb P^{-1}_T}(1)
\]
is regarded as the universal line quotient.  Thus the local projective
space associated to a trivial bundle of rank \(d\geq0\) is
\(\mathbb P^{d-1}_T\), including the rank-zero case.
\end{convention}

\begin{construction}[Projective bundle]
\label{def:spectral-projective-bundle}
Let \(E\) be a finite locally free quasi-coherent
\(\mathcal O_S\)-module.  Choose a Zariski cover
\[
    \{U_a\longrightarrow S\}_{a\in A}
\]
such that, on each \(U_a\), the rank of \(E\) is constant and there is an
equivalence
\[
    E|_{U_a}
    \simeq
    \mathcal O_{U_a}^{\oplus d_a},
    \qquad
    d_a\geq0.
\]
On an overlap
\[
    U_{ab}:=U_a\times_SU_b,
\]
the two trivializations differ by an automorphism of the corresponding
trivial finite locally free module.  Projectivization carries this
automorphism to an automorphism
\[
    \mathbb P^{d_a-1}_{U_{ab}}
    \xrightarrow{\simeq}
    \mathbb P^{d_b-1}_{U_{ab}}.
\]
On every nonempty overlap the ranks agree, and the induced projective
automorphisms satisfy the cocycle condition on triple overlaps.

Spectral Zariski gluing therefore produces a spectral \(S\)-scheme
\[
    \pi:
    \mathbb P_S(E)
    \longrightarrow
    S
\]
together with canonical equivalences
\[
    \mathbb P_S(E)\times_SU_a
    \simeq
    \mathbb P^{d_a-1}_{U_a}.
\]

When \(d_a>0\), the standard quotient
\[
    \mathcal O_{\mathbb P^{d_a-1}_{U_a}}^{\oplus d_a}
    \twoheadrightarrow
    \mathcal O_{\mathbb P^{d_a-1}_{U_a}}(1)
\]
is equivariant for change of trivialization.  When \(d_a=0\), the
rank-zero convention supplies the unique quotient on the empty chart.
These local quotients therefore glue to an invertible module
\[
    \mathcal O_{\mathbb P_S(E)}(1)
\]
and a universal quotient
\[
    \pi^*E
    \twoheadrightarrow
    \mathcal O_{\mathbb P_S(E)}(1).
\]
The resulting spectral scheme and universal quotient are independent of the
chosen trivializing cover through a contractible space of equivalences.
\end{construction}

\begin{proposition}[Representing property and local structure]
\label{prop:projective-bundle-local}
Let \(E\) be finite locally free on \(S\).  For every spectral
\(S\)-scheme
\[
    f:T\longrightarrow S,
\]
composition with the universal quotient induces an equivalence between
\[
    \Map_S
    \bigl(
      T,
      \mathbb P_S(E)
    \bigr)
\]
and the space of line quotients
\[
    f^*E
    \twoheadrightarrow
    L
\]
in the sense of
Definition~\ref{def:spectral-line-quotient}.

In particular, on every open spectral subscheme \(U\subseteq S\) for which
\[
    E|_U
    \simeq
    \mathcal O_U^{\oplus d},
    \qquad
    d\geq0,
\]
there is a canonical equivalence
\[
    \mathbb P_S(E)\times_SU
    \simeq
    \mathbb P^{d-1}_U
\]
compatible with the universal quotient line bundles.
\end{proposition}

\begin{proof}
For a spectral \(S\)-scheme
\[
    f:T\longrightarrow S,
\]
let
\[
    \mathcal Q_E(T)
\]
denote the space of pairs
\[
    (L,q)
\]
where \(L\) is finite locally free of rank one on \(T\) and
\[
    q:f^*E\longrightarrow L
\]
is surjective on \(\pi_0\).  Equivalences in \(\mathcal Q_E(T)\) are
equivalences of line modules compatible with the quotient maps.

The assignment
\[
    T\longmapsto\mathcal Q_E(T)
\]
is a Zariski sheaf of spaces.  Indeed, finite locally free modules, morphisms
between them, equivalences, and surjectivity on \(\pi_0\) all satisfy
Zariski descent.

Suppose first that
\[
    E\simeq0.
\]
If \(T\) is nonempty, there is no line quotient
\[
    0\longrightarrow L,
\]
because \(\pi_0(L)\) is locally free of rank one and the zero map cannot be
surjective at any point of \(T\).  For \(T=\varnothing\), the space of line
quotients is contractible.  Hence
\[
    \mathcal Q_E(T)
    \simeq
    \Map_S
    \bigl(
      T,
      \mathbb P^{-1}_S
    \bigr).
\]
Thus the rank-zero line-quotient functor is represented by the empty
projective space.

We now treat the case
\[
    E\simeq
    \mathcal O_S^{\oplus(r+1)},
    \qquad
    r\geq0.
\]
Let
\[
    e_0,\ldots,e_r
\]
be the standard basis.  For \(0\leq a\leq r\), define
\[
    \mathcal Q_{E,a}(T)
    \subseteq
    \mathcal Q_E(T)
\]
to be the full subspace of line quotients
\[
    q:
    \mathcal O_T^{\oplus(r+1)}
    \longrightarrow
    L
\]
for which the component
\[
    q(e_a):
    \mathcal O_T
    \longrightarrow
    L
\]
is an equivalence.

The equivalence \(q(e_a)\) canonically trivializes \(L\).  Under this
trivialization, \(q\) is uniquely represented by a tuple
\[
    (x_0,\ldots,x_r)
\]
of endomorphisms of \(\mathcal O_T\) with
\[
    x_a=1.
\]
Thus the remaining \(r\) components define a point of brave new affine
\(r\)-space.  This construction gives a natural equivalence of spaces
\[
    \mathcal Q_{E,a}(T)
    \simeq
    \Map_S
    \bigl(
      T,
      \mathbb A^r_S
    \bigr).
\]
Hence
\[
    \mathcal Q_{E,a}
\]
is represented by the standard affine chart
\[
    U_a\simeq\mathbb A^r_S.
\]

On the intersection
\[
    \mathcal Q_{E,a}\cap\mathcal Q_{E,b},
\]
both
\[
    q(e_a)
    \qquad\text{and}\qquad
    q(e_b)
\]
are equivalences.  Under the \(a\)-th trivialization of \(L\), this is the
principal-open condition that the \(b\)-th coordinate be invertible.  Passing
from the \(a\)-th trivialization to the \(b\)-th trivialization divides all
coordinates by the \(b\)-th coordinate.  Consequently, the transition
equivalence on
\[
    \mathcal Q_{E,a}\cap\mathcal Q_{E,b}
\]
is exactly the usual homogeneous-coordinate transition equivalence between
the standard affine charts of
\[
    \mathbb P^r_S.
\]

The subfunctors \(\mathcal Q_{E,a}\) form a Zariski cover of
\(\mathcal Q_E\).  Indeed, surjectivity of
\[
    \pi_0(q):
    \pi_0(\mathcal O_T)^{\oplus(r+1)}
    \longrightarrow
    \pi_0(L)
\]
implies that Zariski-locally on \(T\), at least one component
\[
    \pi_0(q(e_a))
\]
generates the rank-one module \(\pi_0(L)\).  On that open, \(q(e_a)\) is an
equivalence of rank-one finite locally free spectral modules.

It follows that the sheaf of spaces \(\mathcal Q_E\) is obtained by gluing
the affine charts
\[
    \mathbb A^r_S
\]
along precisely the principal-open transition equivalences used in the
construction of
\[
    \mathbb P^r_S.
\]
Therefore there is a natural equivalence
\[
    \Map_S
    \bigl(
      T,
      \mathbb P^r_S
    \bigr)
    \simeq
    \mathcal Q_E(T)
\]
of spaces, not merely a bijection on connected components.

Now let \(E\) be arbitrary finite locally free.  Choose a Zariski cover
\[
    \{U_a\longrightarrow S\}
\]
on which \(E\) is trivial of constant rank.  Both functors
\[
    T\longmapsto
    \Map_S
    \bigl(
      T,
      \mathbb P_S(E)
    \bigr)
\]
and
\[
    T\longmapsto
    \mathcal Q_E(T)
\]
satisfy Zariski descent in \(T\).  On the pullback of each \(U_a\), the
preceding trivial-bundle calculation identifies them.

The local equivalences are natural under automorphisms of the trivial
bundle.  Hence they are compatible on overlaps with the transition
automorphisms used to construct
\[
    \mathbb P_S(E).
\]
Zariski descent therefore gives a natural equivalence
\[
    \Map_S
    \bigl(
      T,
      \mathbb P_S(E)
    \bigr)
    \simeq
    \mathcal Q_E(T).
\]

Under this equivalence, the universal quotient
\[
    \pi^*E
    \twoheadrightarrow
    \mathcal O_{\mathbb P_S(E)}(1)
\]
corresponds to the identity map of
\[
    \mathbb P_S(E).
\]

Finally, if
\[
    E|_U
    \simeq
    \mathcal O_U^{\oplus d},
    \qquad
    d\geq0,
\]
then both
\[
    \mathbb P_S(E)\times_SU
\]
and
\[
    \mathbb P^{d-1}_U
\]
represent the same line-quotient functor on spectral \(U\)-schemes.
Yoneda gives the asserted canonical equivalence, including compatibility
with the universal quotient line bundles.
\end{proof}

\begin{proposition}[Geometry of projective bundles]
\label{prop:projective-bundle-geometry}
Let \(E\) be finite locally free on \(S\).  The structural morphism
\[
    \pi:\mathbb P_S(E)\longrightarrow S
\]
is smooth, proper, and of finite presentation.

For every morphism
\[
    f:T\longrightarrow S,
\]
there is a canonical equivalence
\[
    \mathbb P_S(E)\times_ST
    \simeq
    \mathbb P_T(f^*E),
\]
compatible with the universal quotient line bundles.  Moreover,
\[
    \bigl(\mathbb P_S(E)\bigr)_{\mathrm{cl}}
    \simeq
    \mathbb P_{S_{\mathrm{cl}}}
    \bigl(
      \iota_S^*E
    \bigr).
\]
\end{proposition}

\begin{proof}
The assertions concerning smoothness and finite presentation are local on
\(S\) for the Zariski topology.  After restricting to an open on which
\[
    E\simeq\mathcal O_S^{\oplus d},
    \qquad
    d\geq0,
\]
the structural morphism becomes
\[
    \mathbb P^{d-1}_S\longrightarrow S.
\]
For \(d>0\), this morphism is smooth and of finite presentation by
Proposition~\ref{prop:projective-space-smooth-proper}.  For \(d=0\), it is
the empty morphism, which is smooth and of finite presentation.

The classical truncation is likewise local on \(S\).  On a positive-rank
trivializing open, Proposition~\ref{prop:projective-space-truncation} gives
\[
    (\mathbb P^{d-1}_S)_{\mathrm{cl}}
    \simeq
    \mathbb P^{d-1}_{S_{\mathrm{cl}}}.
\]
On a rank-zero open, both sides are empty.
The local equivalences are compatible with the transition automorphisms, so
they descend to
\[
    \bigl(\mathbb P_S(E)\bigr)_{\mathrm{cl}}
    \simeq
    \mathbb P_{S_{\mathrm{cl}}}
    \bigl(
      \iota_S^*E
    \bigr).
\]
The right-hand side is proper over \(S_{\mathrm{cl}}\).  Since \(\pi\) is of
finite presentation and properness is detected on classical truncations,
\(\pi\) is proper.

For base change, both
\[
    \mathbb P_S(E)\times_ST
\]
and
\[
    \mathbb P_T(f^*E)
\]
represent the functor assigning to a spectral \(T\)-scheme
\[
    h:V\longrightarrow T
\]
the space of invertible quotients
\[
    h^*f^*E\twoheadrightarrow L.
\]
The representing property therefore gives a canonical equivalence between
them.  Naturality and compatibility with the universal quotient follow from
Yoneda.
\end{proof}

\subsection{Projective and proper morphisms}

\begin{definition}[Projective presentation and projective morphism]
\label{def:spectral-projective-morphism}
Let
\[
    f:X\longrightarrow S
\]
be a morphism of spectral schemes.

A projective presentation of \(f\) of length \(m\geq0\) consists of a
factorization
\[
    X=X_m
    \xrightarrow{f_m}
    X_{m-1}
    \xrightarrow{f_{m-1}}
    \cdots
    \xrightarrow{f_2}
    X_1
    \xrightarrow{f_1}
    X_0=S
\]
together with, for every \(1\leq a\leq m\),

\begin{enumerate}
    \item a finite locally free module \(E_a\) on \(X_{a-1}\);

    \item a closed immersion of finite presentation
    \[
        i_a:
        X_a
        \hookrightarrow
        \mathbb P_{X_{a-1}}(E_a);
    \]

    \item an equivalence
    \[
        f_a
        \simeq
        \pi_a i_a,
    \]
    where
    \[
        \pi_a:
        \mathbb P_{X_{a-1}}(E_a)
        \longrightarrow
        X_{a-1}
    \]
    is the structural morphism.
\end{enumerate}

A presentation of length zero is the identity morphism.  A presentation of
length one is called an elementary projective presentation.

The morphism \(f\) is projective if it admits a finite projective
presentation.
\end{definition}

\begin{proposition}[Permanence of projective morphisms]
\label{prop:projective-permanence}
Projective morphisms are stable under equivalence, arbitrary base change,
products over a base, composition, and composition with closed immersions of
finite presentation.

Every projective morphism is proper and of finite presentation.  Classical
truncation carries projective morphisms to classical projective morphisms.
\end{proposition}

\begin{proof}
Stability under equivalence is immediate.

Let
\[
    X=X_m
    \longrightarrow
    X_{m-1}
    \longrightarrow
    \cdots
    \longrightarrow
    X_0=S
\]
be a projective presentation, and let
\[
    T\longrightarrow S
\]
be arbitrary.  Pulling back every term, vector bundle, projective bundle, and
closed immersion gives a projective presentation
\[
    X\times_ST
    \longrightarrow
    T.
\]
Thus projective morphisms are stable under arbitrary base change.

If
\[
    X\xrightarrow{f}Y
    \qquad\text{and}\qquad
    Y\xrightarrow{g}S
\]
admit projective presentations, concatenating the two presentations gives a
projective presentation of
\[
    gf:X\longrightarrow S.
\]
Hence projective morphisms are stable under composition.

Let
\[
    X_1\longrightarrow S,
    \qquad
    X_2\longrightarrow S
\]
be projective.  The projection
\[
    X_1\times_SX_2
    \longrightarrow
    X_2
\]
is the base change of \(X_1\to S\), hence is projective.  Composing it with
the projective morphism \(X_2\to S\) shows that
\[
    X_1\times_SX_2
    \longrightarrow
    S
\]
is projective.

A closed immersion of finite presentation
\[
    i:Z\hookrightarrow Y
\]
is elementary projective: use the canonical equivalence
\[
    \mathbb P_Y(\mathcal O_Y)
    \simeq
    Y
\]
and regard \(i\) as a closed immersion into this rank-one projective bundle.
Consequently, composition with a closed immersion of finite presentation
preserves projectivity.

Every elementary projective morphism is the composite of a closed immersion
of finite presentation with a projective-bundle projection.  Both morphisms
are proper and of finite presentation.  Proper morphisms and morphisms of
finite presentation are stable under finite composition, so every projective
morphism is proper and of finite presentation.

Finally, classical truncation preserves finite locally free modules,
projective bundles, closed immersions of finite presentation, and
composition.  Truncating a finite spectral projective presentation therefore
gives a finite classical projective presentation.  Its composite is a
classical projective morphism.
\end{proof}

\begin{proposition}[Permanence of proper morphisms]
\label{prop:proper-permanence}
Morphisms in \(\mathcal P\) are stable under equivalence, composition,
arbitrary base change, and products over a common base.

Every closed immersion belongs to \(\mathcal P\).  Every morphism in
\(\mathcal P\) is separated: its diagonal is a closed immersion.
\end{proposition}

\begin{proof}
Classical truncation preserves fibre products, finite products, and
composition.  The asserted permanence properties therefore follow from the
corresponding classical permanence properties of proper morphisms.

Every closed immersion belongs to \(\mathcal P\), because its classical
truncation is a closed immersion and hence proper.  If
\[
    p:Y\longrightarrow X
\]
belongs to \(\mathcal P\), then
\[
    \Delta_{p_{\mathrm{cl}}}:
    Y_{\mathrm{cl}}
    \longrightarrow
    Y_{\mathrm{cl}}\times_{X_{\mathrm{cl}}}Y_{\mathrm{cl}}
\]
is a closed immersion.  Classical truncation preserves fibre products, so
\[
    (\Delta_p)_{\mathrm{cl}}
    \simeq
    \Delta_{p_{\mathrm{cl}}}.
\]
Closed immersions of spectral schemes are detected on classical
truncations.  Therefore
\[
    \Delta_p:
    Y\longrightarrow Y\times_XY
\]
is a closed immersion, and \(p\) is separated.
\end{proof}

\section{Stable purity transformations}
\label{sec:stable-purity-transformations}

\subsection{Purity for smooth closed pairs}

\begin{definition}[Smooth closed pair]
\label{def:stable-smooth-closed-pair}
A smooth closed pair over \(S\) is a closed immersion
\[
    i:Z\hookrightarrow X
\]
such that both structural morphisms
\[
    X\longrightarrow S,
    \qquad
    Z\longrightarrow S
\]
are smooth and of finite presentation.
\end{definition}

\begin{construction}[Relative closed-purity transformation]
\label{constr:closed-purity-transformation}
Let
\[
    i:Z\hookrightarrow X
\]
be a smooth closed pair over \(S\), and write
\[
    x:X\longrightarrow S,
    \qquad
    z:Z\longrightarrow S
\]
for the structural morphisms, so that
\[
    z=xi.
\]

Stable comparison and its compatibility with pullback and Thom twists give
a canonical equivalence
\[
    \Theta_Z^{\mathrm{st}}
    \Sigma^{-\nu_i}z^*
    \bigl(
      \Theta_S^{\mathrm{st}}
    \bigr)^{-1}
    \simeq
    \Sigma^{-\nu_{i_{\mathrm{cl}}}}
    z_{\mathrm{cl}}^*.
\]
Compatibility with closed exceptional inverse image gives a canonical
equivalence
\[
    \Theta_Z^{\mathrm{st}}
    i^!x^*
    \bigl(
      \Theta_S^{\mathrm{st}}
    \bigr)^{-1}
    \simeq
    i_{\mathrm{cl}}^!x_{\mathrm{cl}}^*.
\]

Let
\[
    \mathfrak p^{\mathrm{cl}}_{i/S}:
    \Sigma^{-\nu_{i_{\mathrm{cl}}}}
    z_{\mathrm{cl}}^*
    \longrightarrow
    i_{\mathrm{cl}}^!x_{\mathrm{cl}}^*
\]
be the canonical relative-purity transformation in ordinary stable motivic
homotopy theory.

Conjugation by the comparison equivalences induces an equivalence of functor
categories
\[
\begin{aligned}
    \Fun
    \bigl(
      \SH(S),
      \SH(Z)
    \bigr)
    &\xrightarrow{\simeq}
    \Fun
    \bigl(
      \SH_{\mathrm{cl}}(S_{\mathrm{cl}}),
      \SH_{\mathrm{cl}}(Z_{\mathrm{cl}})
    \bigr), \\
    F
    &\longmapsto
    \Theta_Z^{\mathrm{st}}
    F
    \bigl(
      \Theta_S^{\mathrm{st}}
    \bigr)^{-1}.
\end{aligned}
\]
Define
\[
    \mathfrak p_{i/S}:
    \Sigma^{-\nu_i}z^*
    \longrightarrow
    i^!x^*
\]
to be the unique transformation corresponding to
\(\mathfrak p^{\mathrm{cl}}_{i/S}\) under this equivalence.

For the zero section
\[
    s_E:X\hookrightarrow\mathbb V_X(E),
\]
the transformation is normalized by the ordinary inverse Thom equivalence.
Compatibility of stable comparison with Thom spaces identifies this with the
inverse Thom equivalence associated to \(E\) in \(\SH(X)\).
\end{construction}

\begin{theorem}[Relative purity for smooth closed pairs]
\label{thm:closed-purity-smooth-pair}
Let
\[
    i:Z\hookrightarrow X
\]
be a smooth closed pair over \(S\), with structural morphisms
\[
    x:X\longrightarrow S,
    \qquad
    z:Z\longrightarrow S.
\]
The relative purity transformation is an equivalence:
\[
    \mathfrak p_{i/S}:
    \Sigma^{-\nu_i}z^*
    \xrightarrow{\simeq}
    i^!x^*.
\]
It is compatible with arbitrary base change in \(S\), composition of smooth
closed pairs, and virtual Thom twists pulled back from the common base \(S\).
\end{theorem}

\begin{proof}
By Construction~\ref{constr:closed-purity-transformation}, conjugation by
stable classical comparison carries
\[
    \mathfrak p_{i/S}
\]
to the ordinary relative-purity transformation
\[
    \mathfrak p^{\mathrm{cl}}_{i/S}:
    \Sigma^{-\nu_{i_{\mathrm{cl}}}}
    z_{\mathrm{cl}}^*
    \longrightarrow
    i_{\mathrm{cl}}^!x_{\mathrm{cl}}^*.
\]
Since
\[
    i_{\mathrm{cl}}:
    Z_{\mathrm{cl}}
    \hookrightarrow
    X_{\mathrm{cl}}
\]
is a closed immersion between schemes smooth over \(S_{\mathrm{cl}}\),
ordinary relative purity identifies
\(\mathfrak p^{\mathrm{cl}}_{i/S}\) as an equivalence.

The comparison functor
\[
    \Theta_Z^{\mathrm{st}}
\]
is an equivalence and hence conservative.  Therefore
\(\mathfrak p_{i/S}\) is an equivalence.

Ordinary relative purity is compatible with arbitrary base change,
composition of smooth closed pairs, and Thom twists pulled back from the
common base.  Stable comparison is
compatible with each of the operations occurring in those compatibility
diagrams.  Conjugation by comparison therefore transports the ordinary
compatibility diagrams to the asserted spectral compatibility diagrams.
\end{proof}

\begin{remark}[Scope of relative closed purity]
\label{rem:closed-purity-scope}
Theorem~\ref{thm:closed-purity-smooth-pair} asserts purity for objects pulled
back from the common smooth base \(S\).  It does not assert an equivalence
\[
    \Sigma^{-\nu_i}i^*
    \simeq
    i^!
\]
on all of \(\SH(X)\).  That stronger assertion is an absolute-purity
statement and is not required for the diagonal construction below.
\end{remark}

\subsection{Diagonal purity}

Let
\[
    p:X\longrightarrow S
\]
be smooth, separated, and of finite presentation.  Write
\[
    X^{(2)}
    :=
    X\times_SX,
\]
and let
\[
    \pi_1,\pi_2:X^{(2)}\longrightarrow X
\]
be the projections.

\begin{proposition}[The smooth diagonal]
\label{prop:smooth-diagonal-purity}
Let
\[
    p:X\longrightarrow S
\]
be smooth, separated, and of finite presentation.  The diagonal
\[
    \Delta_p:
    X\hookrightarrow X^{(2)}
    :=
    X\times_SX
\]
is a quasi-smooth closed immersion.

Relative to the first projection
\[
    \pi_1:X^{(2)}\longrightarrow X,
\]
the diagram
\[
    X\hookrightarrow X^{(2)}
\]
is a smooth closed pair over \(X\).  Its conormal module is canonically
equivalent to
\[
    N^*_{\Delta_p}
    \simeq
    \Omega_{X/S}.
\]
Consequently, relative purity gives a canonical equivalence of functors
\[
    \Sigma^{-T_p}
    \xrightarrow{\simeq}
    \Delta_p^!\pi_1^*
\]
from \(\SH(X)\) to \(\SH(X)\).
\end{proposition}

\begin{proof}
Since \(p\) is separated, its diagonal
\[
    \Delta_p:X\longrightarrow X^{(2)}
\]
is a closed immersion.

Consider the composite
\[
    X
    \xrightarrow{\Delta_p}
    X^{(2)}
    \xrightarrow{\pi_1}
    X.
\]
Since
\[
    \pi_1\Delta_p=\id_X,
\]
the cotangent transitivity triangle gives
\[
    L_{X/X^{(2)}}[-1]
    \simeq
    \Delta_p^*L_{X^{(2)}/X}.
\]
The morphism \(\pi_1\) is the base change of \(p\), so
\[
    L_{X^{(2)}/X}
    \simeq
    \pi_2^*L_{X/S}.
\]
Pulling back along \(\Delta_p\) gives
\[
    N^*_{\Delta_p}
    =
    L_{X/X^{(2)}}[-1]
    \simeq
    L_{X/S}
    \simeq
    \Omega_{X/S}.
\]
The module \(\Omega_{X/S}\) is finite locally free, so \(\Delta_p\) is
quasi-smooth.

The morphism
\[
    \pi_1:X^{(2)}\longrightarrow X
\]
is smooth of finite presentation, being a base change of \(p\), and
\[
    X\longrightarrow X
\]
is smooth.  Thus \(\Delta_p\) is a smooth closed pair over \(X\).
Applying Theorem~\ref{thm:closed-purity-smooth-pair} with base \(X\) gives
\[
    \Sigma^{-T_p}
    \simeq
    \Delta_p^!\pi_1^*.
\]
\end{proof}

\begin{corollary}[Coherent diagonal Thom equivalence]
\label{cor:coherent-diagonal-thom}
Let
\[
    p:X\longrightarrow S
\]
be smooth, separated, and of finite presentation.  Write
\[
    X^{(2)}:=X\times_SX
\]
with second projection
\[
    \pi_2:X^{(2)}\longrightarrow X
\]
and diagonal
\[
    \Delta_p:X\hookrightarrow X^{(2)}.
\]

There is a canonical equivalence
\[
    \vartheta_p:
    \pi_{2\sharp}\Delta_{p*}\mathbf1_X
    \xrightarrow{\simeq}
    \Sigma^{T_p}\mathbf1_X.
\]
It is compatible with arbitrary base change and with composition of smooth
separated morphisms of finite presentation.
\end{corollary}

\begin{proof}
View
\[
    \Delta_p:
    X\hookrightarrow X^{(2)}
\]
as a smooth closed pair over \(X\) through the second projection
\[
    \pi_2:X^{(2)}\longrightarrow X.
\]
The structural morphism of the closed subscheme is the identity of \(X\),
and its conormal class is
\[
    \nu_{\Delta_p}
    =
    T_p.
\]
Theorem~\ref{thm:closed-purity-smooth-pair} therefore gives a canonical
equivalence of functors
\[
    \Sigma^{-T_p}
    \xrightarrow{\simeq}
    \Delta_p^!\pi_2^*
\]
from \(\SH(X)\) to itself.

Let \(A\in\SH(X)\).  Using the adjunctions
\[
    \pi_{2\sharp}\dashv\pi_2^*,
    \qquad
    \Delta_{p*}\dashv\Delta_p^!,
\]
and relative purity, there are natural equivalences
\[
\begin{aligned}
    \Map_X
    \bigl(
      \pi_{2\sharp}\Delta_{p*}\mathbf1_X,
      A
    \bigr)
    &\simeq
    \Map_{X^{(2)}}
    \bigl(
      \Delta_{p*}\mathbf1_X,
      \pi_2^*A
    \bigr) \\
    &\simeq
    \Map_X
    \bigl(
      \mathbf1_X,
      \Delta_p^!\pi_2^*A
    \bigr) \\
    &\simeq
    \Map_X
    \bigl(
      \mathbf1_X,
      \Sigma^{-T_p}A
    \bigr) \\
    &\simeq
    \Map_X
    \bigl(
      \Sigma^{T_p}\mathbf1_X,
      A
    \bigr).
\end{aligned}
\]
By the covariant Yoneda lemma, this natural equivalence of mapping functors
corresponds to a canonical equivalence
\[
    \widetilde\vartheta_p:
    \Sigma^{T_p}\mathbf1_X
    \xrightarrow{\simeq}
    \pi_{2\sharp}\Delta_{p*}\mathbf1_X.
\]
Define
\[
    \vartheta_p:
    \pi_{2\sharp}\Delta_{p*}\mathbf1_X
    \xrightarrow{\simeq}
    \Sigma^{T_p}\mathbf1_X
\]
to be the inverse of \(\widetilde\vartheta_p\).

The relative-purity transformation of
Theorem~\ref{thm:closed-purity-smooth-pair}
is compatible with arbitrary base change.  The adjunction equivalences used
above and the covariant Yoneda correspondence are likewise natural under
base change.  Consequently, both
\(\widetilde\vartheta_p\) and its inverse \(\vartheta_p\) are compatible
with arbitrary base change.

For composition, let
\[
    X\xrightarrow{p}Y\xrightarrow{q}S
\]
be smooth separated morphisms of finite presentation.  Write
\[
    \kappa:
    X\times_YX
    \longrightarrow
    X\times_SX
\]
for the morphism induced by the two projections to \(X\).  It is the
pullback of
\[
    \Delta_q:
    Y\longrightarrow Y\times_SY
\]
along
\[
    p\times p:
    X\times_SX
    \longrightarrow
    Y\times_SY.
\]
There is a canonical factorization
\[
    \Delta_{qp}
    =
    \kappa\Delta_p.
\]
Both factors are smooth closed pairs over \(X\), with respect to the second
projection.  Compatibility of relative purity with composition identifies
the normal Thom class of the composite with
\[
\begin{aligned}
    \nu_{\Delta_{qp}}
    &\simeq
    \nu_{\Delta_p}
    +
    \Delta_p^*\nu_\kappa \\
    &\simeq
    T_p+p^*T_q.
\end{aligned}
\]
Under the cotangent identifications, this is the path induced by the
transitivity sequence
\[
    p^*\Omega_{Y/S}
    \longrightarrow
    \Omega_{X/S}
    \longrightarrow
    \Omega_{X/Y}.
\]

The composition compatibility of relative purity identifies the
equivalence
\(\widetilde\vartheta_{qp}\) with the pasted equivalence obtained from
\(\widetilde\vartheta_p\) and \(\widetilde\vartheta_q\), under the canonical
identification
\[
    T_{qp}
    \simeq
    T_p+p^*T_q.
\]
Taking inverses therefore identifies the corresponding pasted diagonal Thom
equivalence with \(\vartheta_{qp}\), not merely with an equivalence having
the same source and target.  Associativity of cotangent transitivity and
the composition coherence of relative purity give the corresponding higher
composition coherences.
\end{proof}

\begin{proposition}[Additivity of tangent classes]
\label{prop:tangent-class-composition}
Let
\[
    X\xrightarrow{p}Y\xrightarrow{q}S
\]
be smooth morphisms of finite presentation.  There is a canonical path in
\(K(X)\)
\[
    T_{qp}
    \simeq
    T_p+p^*T_q.
\]
This path is natural under base change and coherently associative for a
triple composite of smooth morphisms.
\end{proposition}

\begin{proof}
The cotangent transitivity sequence is an exact sequence of finite locally
free modules
\[
    p^*\Omega_{Y/S}
    \longrightarrow
    \Omega_{X/S}
    \longrightarrow
    \Omega_{X/Y}.
\]
The coherent Thom additivity theorem of
Chapter~\ref{chap:stable-brave-new-motivic} carries this exact sequence to
the displayed path in vector-bundle \(K\)-theory and hence to the
corresponding equivalence of Thom twists.  Naturality and higher coherence
follow from transitivity for cotangent complexes and the simplicial
coherence of the motivic \(J\)-homomorphism.
\end{proof}

\section{Smooth proper ambidexterity}
\label{sec:smooth-proper-ambidexterity}

\subsection{The ambidexterity transformation}

Let
\[
    p:X\longrightarrow S
\]
be smooth, separated, and of finite presentation.  The two adjunctions already
available are
\[
    p_\sharp\dashv p^*\dashv p_*.
\]
Diagonal purity supplies the twist that compares the two adjoints of
\(p^*\).

\begin{construction}[Ambidexterity transformation]
\label{constr:ambidexterity-transformation}
Let
\[
    p:X\longrightarrow S
\]
be smooth, separated, and of finite presentation.  Write
\[
    X^{(2)}:=X\times_SX
\]
with projections
\[
    \pi_1,\pi_2:X^{(2)}\longrightarrow X
\]
and diagonal
\[
    \Delta_p:X\hookrightarrow X^{(2)}.
\]

Let \(A\in\SH(X)\).  Smooth base change gives a canonical equivalence
\[
    p^*p_\sharp\Sigma^{-T_p}A
    \simeq
    \pi_{2\sharp}\pi_1^*\Sigma^{-T_p}A.
\]
Apply \(\pi_{2\sharp}\) to the unit of
\[
    \Delta_p^*\dashv\Delta_{p*}:
\]
\[
    \pi_1^*\Sigma^{-T_p}A
    \longrightarrow
    \Delta_{p*}\Delta_p^*
    \pi_1^*\Sigma^{-T_p}A.
\]
Since
\[
    \pi_1\Delta_p
    =
    \pi_2\Delta_p
    =
    \id_X,
\]
the target is canonically
\[
    \Delta_{p*}\Sigma^{-T_p}A.
\]

The closed projection formula gives
\[
    \Delta_{p*}\Sigma^{-T_p}A
    \simeq
    \Delta_{p*}\mathbf1_X
    \otimes
    \pi_2^*\Sigma^{-T_p}A.
\]
Applying the smooth projection formula for \(\pi_2\) gives
\[
\begin{aligned}
    \pi_{2\sharp}
    \Delta_{p*}\Sigma^{-T_p}A
    &\simeq
    \pi_{2\sharp}
    \bigl(
      \Delta_{p*}\mathbf1_X
      \otimes
      \pi_2^*\Sigma^{-T_p}A
    \bigr) \\
    &\simeq
    \pi_{2\sharp}\Delta_{p*}\mathbf1_X
    \otimes
    \Sigma^{-T_p}A.
\end{aligned}
\]
The coherent diagonal Thom equivalence of
Corollary~\ref{cor:coherent-diagonal-thom}
identifies this object with
\[
\begin{aligned}
    \Sigma^{T_p}\mathbf1_X
    \otimes
    \Sigma^{-T_p}A
    &\simeq
    A.
\end{aligned}
\]

The resulting composite defines a canonical natural transformation
\[
    \epsilon_p:
    p^*p_\sharp\Sigma^{-T_p}
    \longrightarrow
    \id_{\SH(X)}.
\]
Its right adjoint mate under
\[
    p^*\dashv p_*
\]
is the ambidexterity transformation
\[
    \alpha_p:
    p_\sharp\Sigma^{-T_p}
    \longrightarrow
    p_*.
\]
\end{construction}

\begin{proposition}[Formal properties of ambidexterity]
\label{prop:ambidexterity-formal-properties}
Let
\[
    p:X\longrightarrow S
\]
be smooth, separated, and of finite presentation.

The transformations \(\alpha_p\) satisfy the following properties.

\begin{enumerate}
    \item If \(p=\id_X\), then \(\alpha_p\) is the identity.

    \item If \(p\) is \'{e}tale, then \(T_p=0\), and
    \[
        \alpha_p:
        p_\sharp
        \longrightarrow
        p_*
    \]
    is the untwisted ambidexterity transformation.

    \item For every Cartesian square
    \[
    \begin{tikzcd}[column sep=large,row sep=large]
        X_T \arrow[r,"g'"] \arrow[d,"p'"']
            & X \arrow[d,"p"] \\
        T \arrow[r,"g"'] & S,
    \end{tikzcd}
    \]
    the diagram
    \[
    \begin{tikzcd}[column sep=huge,row sep=large]
        g^*p_\sharp\Sigma^{-T_p}
            \arrow[r,"g^*\alpha_p"]
            \arrow[d,"\simeq"']
            &
        g^*p_*
            \arrow[d,"\operatorname{Ex}^{*}_{*}"]
        \\
        p'_\sharp\Sigma^{-T_{p'}}(g')^*
            \arrow[r,"\alpha_{p'}(g')^*"']
            &
        p'_*(g')^*
    \end{tikzcd}
    \]
    commutes.  The left vertical equivalence is smooth base change together
    with
    \[
        (g')^*T_p\simeq T_{p'},
    \]
    and the right vertical map is the canonical Beck--Chevalley
    transformation.

    \item Let
    \[
        X\xrightarrow{p}Y\xrightarrow{q}S
    \]
    be smooth separated morphisms of finite presentation.  Under the
    canonical equivalences
    \[
        (qp)_\sharp\Sigma^{-T_{qp}}
        \simeq
        q_\sharp\Sigma^{-T_q}
        p_\sharp\Sigma^{-T_p}
    \]
    and
    \[
        (qp)_*
        \simeq
        q_*p_*,
    \]
    the transformation \(\alpha_{qp}\) is the composite
    \[
    \begin{aligned}
        q_\sharp\Sigma^{-T_q}
        p_\sharp\Sigma^{-T_p}
        &\xrightarrow{
          q_\sharp\Sigma^{-T_q}\alpha_p
        }
        q_\sharp\Sigma^{-T_q}p_* \\
        &\xrightarrow{
          \alpha_q p_*
        }
        q_*p_*.
    \end{aligned}
    \]
    The first displayed equivalence uses
    \[
        T_{qp}
        \simeq
        T_p+p^*T_q
    \]
    and the smooth projection formula for \(p_\sharp\).
\end{enumerate}

All four assertions are coherently compatible with identities, iterated base
change, and triple composition.
\end{proposition}

\begin{proof}
For the identity morphism, the fibre product, both projections, the
diagonal, the tangent class, and every adjunction occurring in
Construction~\ref{constr:ambidexterity-transformation}
are identities.  Hence the resulting transformation is the identity.

If \(p\) is \'{e}tale, then
\[
    \Omega_{X/S}\simeq0
\]
and therefore
\[
    T_p=0.
\]
The construction consequently becomes untwisted.

Consider a Cartesian square as in part \(3\).  Every ingredient in the
construction of
\[
    \epsilon_p:
    p^*p_\sharp\Sigma^{-T_p}
    \longrightarrow
    \id_{\SH(X)}
\]
is compatible with arbitrary base change:

\begin{enumerate}
    \item smooth Beck--Chevalley for the two projections;

    \item base change of the diagonal and of its adjunction;

    \item the closed and smooth projection formulas;

    \item pullback of the tangent Thom twist;

    \item the base-change compatibility of the coherent diagonal Thom
    equivalence
    \[
        \vartheta_p:
        \pi_{2\sharp}\Delta_{p*}\mathbf1_X
        \simeq
        \Sigma^{T_p}\mathbf1_X.
    \]
\end{enumerate}

Hence base change carries the counit construction for \(p\) to the counit
construction for \(p'\).  Taking right adjoint mates gives the commutative
square in part \(3\).  No invertibility of the right vertical
Beck--Chevalley transformation is used.

Now let
\[
    X\xrightarrow{p}Y\xrightarrow{q}S
\]
be composable smooth separated morphisms of finite presentation.
Cotangent transitivity gives the coherently associative path
\[
    T_{qp}
    \simeq
    T_p+p^*T_q.
\]
Consequently, the smooth projection formula identifies
\[
    (qp)_\sharp\Sigma^{-T_{qp}}
\]
with
\[
    q_\sharp\Sigma^{-T_q}
    p_\sharp\Sigma^{-T_p}.
\]

The diagonal construction for \(qp\) is obtained by pasting the diagonal
constructions for \(p\) and \(q\).  The composition compatibility of
Corollary~\ref{cor:coherent-diagonal-thom}
identifies the resulting pasted Thom equivalence with
\[
    \vartheta_{qp}.
\]
The pasting laws for smooth base change, projection formulas, units, and
adjoint mates therefore identify the pasted counit with
\(\epsilon_{qp}\).  Taking right adjoint mates gives precisely the composite
displayed in part \(4\).

The higher identity, base-change, and composition coherences follow from the
higher pasting laws for adjunctions and from the coherent base-change and
composition properties of relative purity and the motivic
\(J\)-homomorphism.
\end{proof}

\begin{proposition}[Module-linearity of ambidexterity]
\label{prop:ambidexterity-module-linearity}
Let
\[
    p:X\longrightarrow S
\]
be smooth, separated, and of finite presentation.  For
\[
    A\in\SH(X),
    \qquad
    B\in\SH(S),
\]
the diagram
\[
\begin{tikzcd}[column sep=huge,row sep=large]
    p_\sharp\Sigma^{-T_p}A\otimes_SB
        \arrow[r,"\simeq"]
        \arrow[d,"\alpha_p(A)\otimes\id_B"']
        &
    p_\sharp
    \Sigma^{-T_p}
    \bigl(
      A\otimes_Xp^*B
    \bigr)
        \arrow[d,"\alpha_p(A\otimes p^*B)"]
    \\
    p_*A\otimes_SB
        \arrow[r,"\operatorname{PF}_p"']
        &
    p_*
    \bigl(
      A\otimes_Xp^*B
    \bigr)
\end{tikzcd}
\]
commutes.  Thus
\[
    \alpha_p:
    p_\sharp\Sigma^{-T_p}
    \longrightarrow
    p_*
\]
is a morphism from the strong \(\SH(S)\)-module structure on
\(p_\sharp\Sigma^{-T_p}\) to the canonical lax \(\SH(S)\)-module
structure on \(p_*\).
\end{proposition}

\begin{proof}
Put
\[
    F_p
    :=
    p_\sharp\Sigma^{-T_p}.
\]
The smooth projection formula gives \(F_p\) its strong
\(\SH(S)\)-module structure, while the canonical projection
transformation gives \(p_*\) its lax \(\SH(S)\)-module structure.

It is enough to prove that the transformation
\[
    \epsilon_p:
    p^*F_p
    \longrightarrow
    \id_{\SH(X)}
\]
of Construction~\ref{constr:ambidexterity-transformation} is
\(\SH(S)\)-module-linear.  Under smooth base change,
\[
    p^*F_p(A)
    \simeq
    \pi_{2\sharp}\pi_1^*\Sigma^{-T_p}A.
\]
For \(B\in\SH(S)\), the equality
\[
    p\pi_1
    =
    p\pi_2
\]
gives a canonical equivalence
\[
    \pi_1^*p^*B
    \simeq
    \pi_2^*p^*B.
\]
After this identification, every step in the construction of
\(\epsilon_p(A\otimes p^*B)\) is obtained from the corresponding step for
\(A\) by tensoring with \(\pi_2^*p^*B\): the unit of
\(\Delta_p^*\dashv\Delta_{p*}\), the closed projection formula, the smooth
projection formula for \(\pi_2\), the coherent diagonal Thom equivalence,
and cancellation of the inverse Thom twists.  Consequently, the diagram
\[
\begin{tikzcd}[column sep=huge,row sep=large]
    p^*F_p(A)\otimes_Xp^*B
        \arrow[r,"\simeq"]
        \arrow[d,"\epsilon_p(A)\otimes\id"']
        &
    p^*F_p(A\otimes_Xp^*B)
        \arrow[d,"\epsilon_p(A\otimes p^*B)"]
    \\
    A\otimes_Xp^*B
        \arrow[r,equal]
        &
    A\otimes_Xp^*B
\end{tikzcd}
\]
commutes.

Taking the right adjoint mate under
\[
    p^*\dashv p_*
\]
carries this module-linearity square to the asserted square for
\(\alpha_p\).  The lower horizontal map is precisely the canonical
projection transformation for the right adjoint \(p_*\).
\end{proof}

\begin{theorem}[Smooth proper ambidexterity]
\label{thm:smooth-proper-ambidexterity}
Let
\[
    p:X\longrightarrow S
\]
be smooth, proper, and of finite presentation.  The canonical ambidexterity
transformation is an equivalence:
\[
    \alpha_p:
    p_\sharp\Sigma^{-T_p}
    \xrightarrow{\simeq}
    p_*.
\]
Equivalently,
\[
    p_\sharp
    \xrightarrow{\simeq}
    p_*\Sigma^{T_p}.
\]

The equivalence is \(\SH(S)\)-module-linear, where \(p_*\) carries its
canonical lax \(\SH(S)\)-module structure.  In particular, the canonical
projection transformation for \(p_*\) is an equivalence, so this lax module
structure is strong.

The equivalence is compatible with arbitrary base change and composition.
It is natural in the object of \(\SH(X)\).  Consequently, for every
\[
    \eta\in K(X)
\]
and every \(A\in\SH(X)\), it gives a canonical equivalence
\[
    p_\sharp
    \Sigma^{-T_p}
    \Sigma^\eta A
    \xrightarrow{\simeq}
    p_*\Sigma^\eta A.
\]
If
\[
    \eta=p^*\xi
\]
for some \(\xi\in K(S)\), module-linearity and this projection equivalence
further identify the preceding equivalence with
\[
    p_\sharp
    \Sigma^{-T_p}
    \Sigma^{p^*\xi}A
    \xrightarrow{\simeq}
    \Sigma^\xi p_*A.
\]
\end{theorem}

\begin{proof}
Stable classical comparison carries every ingredient of
Construction~\ref{constr:ambidexterity-transformation} to its ordinary
counterpart: smooth base change, the diagonal adjunction, the
coefficientwise relative closed-purity transformation of
Construction~\ref{constr:closed-purity-transformation} and its diagonal
Thom mate from Corollary~\ref{cor:coherent-diagonal-thom}, the smooth and
closed projection formulas, and the tangent Thom twist.  Consequently, it
carries
\[
    \alpha_p:
    p_\sharp\Sigma^{-T_p}
    \longrightarrow
    p_*
\]
to the ordinary smooth-proper ambidexterity transformation
\[
    (p_{\mathrm{cl}})_\sharp
    \Sigma^{-T_{p_{\mathrm{cl}}}}
    \longrightarrow
    (p_{\mathrm{cl}})_*.
\]
The latter is an equivalence by ordinary motivic smooth-proper
ambidexterity; see, for example,
\cite{HoyoisSixOperations}.  Since
\[
    \Theta_S^{\mathrm{st}}
\]
is conservative, \(\alpha_p\) is an equivalence.

The equivalent form
\[
    p_\sharp
    \simeq
    p_*\Sigma^{T_p}
\]
is obtained by precomposing with the inverse Thom twist.

The base-change and composition compatibilities hold for the canonical
transformation by
Proposition~\ref{prop:ambidexterity-formal-properties}.  Its
\(\SH(S)\)-module-linearity is
Proposition~\ref{prop:ambidexterity-module-linearity}.  These properties
therefore remain valid after \(\alpha_p\) is identified as an equivalence.
Since the source
\[
    p_\sharp\Sigma^{-T_p}
\]
has a strong \(\SH(S)\)-module structure, module-linearity of the
equivalence implies that the canonical lax module structure on \(p_*\) is
strong.  Equivalently, the canonical projection transformation
\[
    p_*A\otimes_SB
    \longrightarrow
    p_*
    \bigl(
      A\otimes_Xp^*B
    \bigr)
\]
is an equivalence.

Naturality of \(\alpha_p\) in
\[
    A\in\SH(X)
\]
allows it to be evaluated on any Thom-twisted object
\[
    \Sigma^\eta A,
    \qquad
    \eta\in K(X),
\]
and gives
\[
    p_\sharp
    \Sigma^{-T_p}
    \Sigma^\eta A
    \xrightarrow{\simeq}
    p_*\Sigma^\eta A.
\]
No movement of an arbitrary source twist through \(p_*\) is asserted here.

If
\[
    \eta=p^*\xi
\]
for some \(\xi\in K(S)\), pullback naturality of the motivic
\(J\)-homomorphism gives
\[
    \Sigma^{p^*\xi}A
    \simeq
    A\otimes_Xp^*\operatorname{th}_S(\xi).
\]
The projection equivalence just established then gives
\[
\begin{aligned}
    p_*\Sigma^{p^*\xi}A
    &\simeq
    p_*
    \bigl(
      A\otimes_Xp^*\operatorname{th}_S(\xi)
    \bigr) \\
    &\simeq
    p_*A\otimes_S\operatorname{th}_S(\xi) \\
    &\simeq
    \Sigma^\xi p_*A.
\end{aligned}
\]
This proves the final displayed compatibility.
\end{proof}

\subsection{Motivic Atiyah duality}

Let \(p:X\to S\) be smooth and proper.  Define
\[
    M_S(X)
    :=
    p_\sharp\mathbf1_X
\]
and
\[
    D_S(X)
    :=
    p_\sharp\Sigma^{-T_p}\mathbf1_X.
\]

\begin{construction}[Coevaluation]
\label{constr:atiyah-coevaluation}
For \(B,C\in\SH(S)\), the smooth projection formula, the adjunctions
\[
    p_\sharp\dashv p^*\dashv p_*,
\]
and smooth proper ambidexterity give natural equivalences
\[
\begin{aligned}
    \Map_S
    \bigl(
      M_S(X)\otimes B,
      C
    \bigr)
    &\simeq
    \Map_S
    \bigl(
      p_\sharp p^*B,
      C
    \bigr) \\
    &\simeq
    \Map_X
    \bigl(
      p^*B,
      p^*C
    \bigr) \\
    &\simeq
    \Map_S
    \bigl(
      B,
      p_*p^*C
    \bigr) \\
    &\simeq
    \Map_S
    \bigl(
      B,
      p_\sharp\Sigma^{-T_p}p^*C
    \bigr) \\
    &\simeq
    \Map_S
    \bigl(
      B,
      D_S(X)\otimes C
    \bigr).
\end{aligned}
\]
Thus
\[
    D_S(X)\otimes-
\]
is right adjoint to
\[
    M_S(X)\otimes-.
\]
The unit of this adjunction, evaluated at \(\mathbf1_S\), is the canonical
coevaluation
\[
    \operatorname{coev}_p:
    \mathbf1_S
    \longrightarrow
    D_S(X)\otimes M_S(X).
\]
Using the symmetry of the tensor product, we also regard it as a map
\[
    \mathbf1_S
    \longrightarrow
    M_S(X)\otimes D_S(X).
\]
\end{construction}

\begin{construction}[Evaluation]
\label{constr:atiyah-evaluation}
In the notation of
Construction~\ref{constr:atiyah-coevaluation}, the counit of the adjunction
\[
    M_S(X)\otimes-
    \dashv
    D_S(X)\otimes-
\]
evaluated at \(\mathbf1_S\) is the canonical evaluation
\[
    \operatorname{ev}_p:
    M_S(X)\otimes D_S(X)
    \longrightarrow
    \mathbf1_S.
\]
After applying the symmetry isomorphism, it may equivalently be written
\[
    \operatorname{ev}_p:
    D_S(X)\otimes M_S(X)
    \longrightarrow
    \mathbf1_S.
\]
\end{construction}

\begin{theorem}[Motivic Atiyah duality]
\label{thm:motivic-atiyah-duality}
Let
\[
    p:X\longrightarrow S
\]
be smooth and proper.  The maps
\[
    \operatorname{coev}_p:
    \mathbf1_S
    \longrightarrow
    M_S(X)\otimes D_S(X)
\]
and
\[
    \operatorname{ev}_p:
    D_S(X)\otimes M_S(X)
    \longrightarrow
    \mathbf1_S
\]
exhibit \(D_S(X)\) as the strong dual of \(M_S(X)\).  Thus
\[
    \bigl(
      p_\sharp\mathbf1_X
    \bigr)^\vee
    \simeq
    p_\sharp\Sigma^{-T_p}\mathbf1_X.
\]
The duality is compatible with arbitrary base change and composition of
smooth proper morphisms.
\end{theorem}

\begin{proof}
Construction~\ref{constr:atiyah-coevaluation} gives an adjunction
\[
    M_S(X)\otimes-
    \dashv
    D_S(X)\otimes-.
\]
In a symmetric monoidal category, an object \(M\) is strongly dualizable with
dual \(D\) precisely when tensoring by \(D\) is right adjoint to tensoring by
\(M\).  The unit and counit of that adjunction are the coevaluation and
evaluation maps, and their triangle identities are the triangle identities
of the adjunction.  Hence \(D_S(X)\) is the strong dual of \(M_S(X)\).

Arbitrary base change preserves the smooth projection formula, the
adjunctions, the tangent Thom class, and the ambidexterity equivalence.
It therefore preserves the tensor adjunction above and its unit and counit.
Compatibility with composition follows in the same way from
Proposition~\ref{prop:tangent-class-composition} and the composition
compatibility of smooth proper ambidexterity.
\end{proof}

\begin{corollary}[\'{E}tale proper ambidexterity]
\label{cor:etale-proper-ambidexterity}
If \(p:X\to S\) is \'{e}tale and proper, then
\[
    p_\sharp
    \simeq
    p_*.
\]
\end{corollary}

\begin{proof}
For an \'{e}tale morphism,
\[
    \Omega_{X/S}\simeq0,
\]
so \(T_p=0\).  Apply
Theorem~\ref{thm:smooth-proper-ambidexterity}.
\end{proof}

\begin{proposition}[Internal Hom and arbitrary direct image]
\label{prop:arbitrary-direct-image-internal-hom}
Let
\[
    f:X\longrightarrow S
\]
be any morphism between quasi-compact quasi-separated spectral schemes.  For
every
\[
    B\in\SH(S),
    \qquad
    A\in\SH(X),
\]
there is a canonical equivalence
\[
    f_*
    \underline{\operatorname{Hom}}_X
    \bigl(
      f^*B,
      A
    \bigr)
    \simeq
    \underline{\operatorname{Hom}}_S
    \bigl(
      B,
      f_*A
    \bigr).
\]
The equivalence is natural in \(A\), \(B\), and \(f\).
\end{proposition}

\begin{proof}
Let \(C\in\SH(S)\).  Using the adjunction \(f^*\dashv f_*\), closedness of
the tensor products, and strong symmetric monoidality of \(f^*\), there are
natural equivalences
\[
\begin{aligned}
    \Map_S
    \Bigl(
      C,\,
      f_*
      \underline{\operatorname{Hom}}_X
      \bigl(
        f^*B,
        A
      \bigr)
    \Bigr)
    &\simeq
    \Map_X
    \Bigl(
      f^*C,\,
      \underline{\operatorname{Hom}}_X
      \bigl(
        f^*B,
        A
      \bigr)
    \Bigr) \\
    &\simeq
    \Map_X
    \bigl(
      f^*C\otimes f^*B,
      A
    \bigr) \\
    &\simeq
    \Map_X
    \bigl(
      f^*(C\otimes B),
      A
    \bigr) \\
    &\simeq
    \Map_S
    \bigl(
      C\otimes B,
      f_*A
    \bigr) \\
    &\simeq
    \Map_S
    \Bigl(
      C,\,
      \underline{\operatorname{Hom}}_S
      \bigl(
        B,
        f_*A
      \bigr)
    \Bigr).
\end{aligned}
\]
Yoneda gives the asserted equivalence.
\end{proof}

\begin{theorem}[Cocontinuity of arbitrary direct image]
\label{thm:arbitrary-direct-image-cocontinuous}
Let
\[
    f:X\longrightarrow S
\]
be any morphism between quasi-compact quasi-separated spectral schemes.  The
functor
\[
    f_*:\SH(X)\longrightarrow\SH(S)
\]
preserves all small colimits.
\end{theorem}

\begin{proof}
Let
\[
    \mathcal G_S
\]
be the family of compact generators constructed in
Chapter~\ref{chap:stable-brave-new-motivic}.  Its members have the form
\[
    \Sigma^{n\mathcal O_S}
    \Sigma_{T,+}^{\infty}\Mot_S(U),
    \qquad
    U\in\Sm^{\mathrm{fp}}_{/S},
    \quad
    n\in\mathbb Z.
\]
Base change of smooth morphisms and pullback compatibility of Thom twists
give
\[
\begin{aligned}
    f^*
    \bigl(
      \Sigma^{n\mathcal O_S}
      \Sigma_{T,+}^{\infty}\Mot_S(U)
    \bigr)
    &\simeq
    \Sigma^{n\mathcal O_X}
    \Sigma_{T,+}^{\infty}
    \Mot_X(U\times_SX).
\end{aligned}
\]
The object on the right is compact in \(\SH(X)\).  By the compact-closure
lemma of Chapter~\ref{chap:stable-brave-new-motivic}, every compact object of
\(\SH(S)\) is a retract of an object obtained from the family
\(\mathcal G_S\) by finite colimits.  The functor \(f^*\) is exact, hence
preserves finite colimits, and it preserves retracts.  Since it carries every
member of \(\mathcal G_S\) to a compact object, it therefore preserves every
compact object.

Let
\[
    A:I\longrightarrow\SH(X)
\]
be a filtered diagram.  For every \(G\in\mathcal G_S\), compactness of
\(G\) and \(f^*G\) gives
\[
\begin{aligned}
    \Map_S
    \bigl(
      G,
      \operatorname*{colim}_{i\in I}f_*A_i
    \bigr)
    &\simeq
    \operatorname*{colim}_{i\in I}
    \Map_S(G,f_*A_i) \\
    &\simeq
    \operatorname*{colim}_{i\in I}
    \Map_X(f^*G,A_i) \\
    &\simeq
    \Map_X
    \bigl(
      f^*G,
      \operatorname*{colim}_{i\in I}A_i
    \bigr) \\
    &\simeq
    \Map_S
    \bigl(
      G,
      f_*\operatorname*{colim}_{i\in I}A_i
    \bigr).
\end{aligned}
\]
The family \(\mathcal G_S\) detects equivalences.  Hence \(f_*\) preserves
filtered colimits.

As a right adjoint between stable \(\infty\)-categories, \(f_*\) is exact
and therefore preserves finite colimits.  An arbitrary coproduct is the
filtered colimit of its finite subcoproducts, so \(f_*\) preserves arbitrary
coproducts.  An exact functor between stable presentable
\(\infty\)-categories that preserves arbitrary coproducts preserves all
small colimits.  Therefore \(f_*\) is colimit-preserving.
\end{proof}

\section{Projective base change and projection}
\label{sec:projective-exchange}

\subsection{Projective bundles}

Consider a Cartesian square
\[
\begin{tikzcd}[column sep=large,row sep=large]
    \mathbb P_T(f^*E) \arrow[r,"g'"] \arrow[d,"\pi'"']
        & \mathbb P_S(E) \arrow[d,"\pi"] \\
    T \arrow[r,"g"'] & S.
\end{tikzcd}
\]
The morphisms \(\pi\) and \(\pi'\) are smooth and proper.

\begin{theorem}[Base change for projective bundles]
\label{thm:projective-bundle-base-change}
The canonical Beck--Chevalley transformation
\[
    g^*\pi_*
    \longrightarrow
    \pi'_*(g')^*
\]
is an equivalence.
\end{theorem}

\begin{proof}
Smooth proper ambidexterity gives
\[
    \pi_*
    \simeq
    \pi_\sharp\Sigma^{-T_\pi}
\]
and
\[
    \pi'_*
    \simeq
    \pi'_\sharp\Sigma^{-T_{\pi'}}.
\]
Base change for cotangent complexes gives
\[
    (g')^*T_\pi
    \simeq
    T_{\pi'}.
\]
Hence stable smooth base change and pullback compatibility of Thom twists
identify
\[
\begin{aligned}
    g^*\pi_*
    &\simeq
    g^*\pi_\sharp\Sigma^{-T_\pi} \\
    &\simeq
    \pi'_\sharp(g')^*\Sigma^{-T_\pi} \\
    &\simeq
    \pi'_\sharp\Sigma^{-T_{\pi'}}(g')^* \\
    &\simeq
    \pi'_*(g')^*.
\end{aligned}
\]
The resulting equivalence is the canonical Beck--Chevalley transformation,
because every step is obtained by taking canonical mates.
\end{proof}




\begin{theorem}[Projection formula for projective bundles]
\label{thm:projective-bundle-projection}
For \(A\in\SH(\mathbb P_S(E))\) and \(B\in\SH(S)\), the canonical
projection morphism
\[
    \pi_*A\otimes_SB
    \longrightarrow
    \pi_*(A\otimes\pi^*B)
\]
is an equivalence.
\end{theorem}

\begin{proof}
The projective-bundle projection
\[
    \pi:\mathbb P_S(E)\longrightarrow S
\]
is smooth and proper.  Hence smooth proper ambidexterity gives equivalences
\[
    \alpha_\pi(A):
    \pi_\sharp\Sigma^{-T_\pi}A
    \xrightarrow{\simeq}
    \pi_*A
\]
and
\[
    \alpha_\pi(A\otimes\pi^*B):
    \pi_\sharp
    \Sigma^{-T_\pi}
    (A\otimes\pi^*B)
    \xrightarrow{\simeq}
    \pi_*(A\otimes\pi^*B).
\]

By
Proposition~\ref{prop:ambidexterity-module-linearity},
there is a commutative square
\[
\begin{tikzcd}[column sep=huge,row sep=large]
    \pi_\sharp\Sigma^{-T_\pi}A\otimes_SB
        \arrow[r,"\simeq"]
        \arrow[d,"\simeq"']
        &
    \pi_\sharp
    \Sigma^{-T_\pi}
    (A\otimes\pi^*B)
        \arrow[d,"\simeq"]
    \\
    \pi_*A\otimes_SB
        \arrow[r,"\operatorname{PF}_\pi"']
        &
    \pi_*(A\otimes\pi^*B).
\end{tikzcd}
\]
The upper horizontal map is the smooth projection equivalence, and both
vertical maps are ambidexterity equivalences.  Therefore the lower
horizontal map, which is the canonical projection transformation, is an
equivalence.
\end{proof}

\subsection{Projective morphisms}

\begin{theorem}[Projective base change]
\label{thm:projective-base-change}
Consider a Cartesian square
\[
\begin{tikzcd}[column sep=large,row sep=large]
    X_T \arrow[r,"g'"] \arrow[d,"f'"']
        & X \arrow[d,"f"] \\
    T \arrow[r,"g"'] & S,
\end{tikzcd}
\]
where \(f\) is projective.  The canonical Beck--Chevalley transformation
\[
    g^*f_*
    \longrightarrow
    f'_*(g')^*
\]
is an equivalence.
\end{theorem}

\begin{proof}
Choose a projective presentation of \(f\), and argue by induction on its
length.

Suppose first that \(f\) admits an elementary projective presentation
\[
    X
    \xrightarrow{i}
    \mathbb P_S(E)
    \xrightarrow{\pi}
    S,
\]
where \(i\) is a closed immersion of finite presentation.  After base change
to \(T\), write
\[
    X_T
    \xrightarrow{i'}
    \mathbb P_T(g^*E)
    \xrightarrow{\pi'}
    T
\]
for the pulled-back presentation, and let
\[
    \widetilde g:
    \mathbb P_T(g^*E)
    \longrightarrow
    \mathbb P_S(E)
\]
be the base-change morphism.

Since
\[
    f_*
    \simeq
    \pi_*i_*,
    \qquad
    f'_*
    \simeq
    \pi'_*i'_*,
\]
stable closed base change and
Theorem~\ref{thm:projective-bundle-base-change} give
\[
\begin{aligned}
    g^*f_*
    &\simeq
    g^*\pi_*i_* \\
    &\simeq
    \pi'_*\widetilde g^*i_* \\
    &\simeq
    \pi'_*i'_*(g')^* \\
    &\simeq
    f'_*(g')^*.
\end{aligned}
\]
By the pasting law for adjoint mates, this composite is the canonical
Beck--Chevalley transformation.

Now suppose that \(f\) has positive presentation length and write
\[
    X\xrightarrow{h}Y\xrightarrow{q}S,
\]
where \(h\) is elementary projective and \(q\) has shorter projective
presentation.  Put
\[
    Y_T:=Y\times_ST
\]
and write
\[
    \widetilde g:Y_T\longrightarrow Y,
    \qquad
    h_T:X_T\longrightarrow Y_T,
    \qquad
    q_T:Y_T\longrightarrow T
\]
for the base-change morphisms.

Using the inductive hypothesis for \(q\) and the elementary case for \(h\),
we obtain
\[
\begin{aligned}
    g^*f_*
    &\simeq
    g^*q_*h_* \\
    &\xrightarrow{\simeq}
    q_{T*}\widetilde g^*h_* \\
    &\xrightarrow{\simeq}
    q_{T*}h_{T*}(g')^* \\
    &\simeq
    f'_*(g')^*.
\end{aligned}
\]
The Beck--Chevalley pasting law identifies this composite with the canonical
base-change transformation for \(f=qh\).  The induction is complete.
\end{proof}

\begin{theorem}[Projective projection formula]
\label{thm:projective-projection-formula}
Let
\[
    f:X\longrightarrow S
\]
be projective.  For
\[
    A\in\SH(X),
    \qquad
    B\in\SH(S),
\]
the canonical projection morphism
\[
    f_*A\otimes_SB
    \longrightarrow
    f_*(A\otimes_Xf^*B)
\]
is an equivalence.
\end{theorem}

\begin{proof}
Choose a projective presentation of \(f\), and argue by induction on its
length.

Suppose first that
\[
    f=\pi i
\]
is elementary projective, where
\[
    i:X\hookrightarrow\mathbb P_S(E)
\]
is a closed immersion and
\[
    \pi:\mathbb P_S(E)\longrightarrow S
\]
is the projective-bundle projection.

The canonical projection transformation for the composite \(f=\pi i\) is,
by the projection-formula pasting law, the composite
\[
\begin{aligned}
    \pi_*i_*A\otimes_SB
    &\longrightarrow
    \pi_*
    \bigl(
      i_*A\otimes\pi^*B
    \bigr) \\
    &\longrightarrow
    \pi_*i_*
    \bigl(
      A\otimes i^*\pi^*B
    \bigr).
\end{aligned}
\]
The first arrow is the projective-bundle projection transformation and is
an equivalence by
Theorem~\ref{thm:projective-bundle-projection}.
The second arrow is induced by the closed projection equivalence
\[
    i_*A\otimes\pi^*B
    \xrightarrow{\simeq}
    i_*
    \bigl(
      A\otimes i^*\pi^*B
    \bigr).
\]
Hence the canonical projection transformation for \(f\) is an equivalence
in the elementary case.

Now write a projective morphism of positive presentation length as
\[
    X\xrightarrow{h}Y\xrightarrow{q}S,
\]
where \(h\) is elementary projective and \(q\) has shorter projective
presentation.  The canonical projection transformation for \(f=qh\) is the
composite
\[
\begin{aligned}
    q_*h_*A\otimes_SB
    &\longrightarrow
    q_*
    \bigl(
      h_*A\otimes_Yq^*B
    \bigr) \\
    &\longrightarrow
    q_*h_*
    \bigl(
      A\otimes_Xh^*q^*B
    \bigr).
\end{aligned}
\]
The first arrow is an equivalence by the induction hypothesis for \(q\).
The second is obtained by applying \(q_*\) to the projection equivalence for
the elementary projective morphism \(h\).  It is therefore an equivalence.

Since
\[
    h^*q^*B
    \simeq
    f^*B,
\]
the target is
\[
    f_*(A\otimes_Xf^*B).
\]
Thus the canonical projection transformation for \(f\) is an equivalence.
\end{proof}

\section{Proper base change and the proper projection formula}
\label{sec:proper-exchange}

\subsection{Canonical exchange transformations}

Consider a Cartesian square
\[
\begin{tikzcd}[column sep=large,row sep=large]
    X_T \arrow[r,"g'"] \arrow[d,"f'"']
        & X \arrow[d,"f"] \\
    T \arrow[r,"g"'] & S.
\end{tikzcd}
\]
The inverse-image composition equivalence gives, by adjunction, a canonical
Beck--Chevalley transformation
\[
    \operatorname{Ex}^{*}_{*}:
    g^*f_*
    \longrightarrow
    f'_*(g')^*.
\]
For any morphism \(f\), lax symmetric monoidality of \(f_*\) gives the
canonical projection morphism
\[
    \operatorname{PF}_f:
    f_*A\otimes_SB
    \longrightarrow
    f_*(A\otimes_Xf^*B).
\]

\begin{lemma}[Comparison of proper transformations]
\label{lem:comparison-proper-transformations}
Suppose that \(f\) is proper.  Under stable classical comparison, the
spectral Beck--Chevalley and projection transformations are carried to the
ordinary transformations for
\[
    f_{\mathrm{cl}}:
    X_{\mathrm{cl}}
    \longrightarrow
    S_{\mathrm{cl}}.
\]
\end{lemma}

\begin{proof}
The transformation \(\operatorname{Ex}^{*}_{*}\) is characterized as the
adjoint mate of the canonical equivalence between the two inverse-image
composites around the Cartesian square.  Stable classical comparison is
compatible with inverse image, direct image, units, counits, and Cartesian
composition.  It therefore carries this mate to the ordinary mate.

The projection transformation is adjoint to
\[
    f^*(f_*A\otimes B)
    \simeq
    f^*f_*A\otimes f^*B
    \longrightarrow
    A\otimes f^*B.
\]
Strong symmetric monoidality of comparison and compatibility with the
counit of \(f^*\dashv f_*\) identify its image with the ordinary projection
transformation.
\end{proof}

\begin{theorem}[Proper base change]
\label{thm:proper-base-change}
For every Cartesian square
\[
\begin{tikzcd}[column sep=large,row sep=large]
    X_T \arrow[r,"g'"] \arrow[d,"f'"']
        & X \arrow[d,"f"] \\
    T \arrow[r,"g"'] & S
\end{tikzcd}
\]
with \(f\) proper, the canonical transformation
\[
    g^*f_*
    \xrightarrow{\simeq}
    f'_*(g')^*
\]
is an equivalence.
\end{theorem}

\begin{proof}
Properness is detected on classical truncations, so
\(f_{\mathrm{cl}}\) is proper.  By
Lemma~\ref{lem:comparison-proper-transformations}, stable classical
comparison carries the displayed transformation to the ordinary proper
base-change transformation.  The latter is an equivalence in ordinary
stable motivic homotopy theory.  Since
\(\Theta_T^{\mathrm{st}}\) is conservative, the spectral transformation is
an equivalence.
\end{proof}

\begin{theorem}[Proper projection formula]
\label{thm:proper-projection-formula}
Let \(f:X\to S\) be proper.  For every
\(A\in\SH(X)\) and \(B\in\SH(S)\), the canonical projection morphism
\[
    f_*A\otimes_SB
    \xrightarrow{\simeq}
    f_*(A\otimes_Xf^*B)
\]
is an equivalence.
\end{theorem}

\begin{proof}
By Lemma~\ref{lem:comparison-proper-transformations}, the image of the
canonical projection morphism under stable classical comparison is the
ordinary proper projection transformation.  The latter is an equivalence.
Conservativity of \(\Theta_S^{\mathrm{st}}\) proves the assertion.
\end{proof}

\begin{remark}[Why the comparison argument is not circular]
\label{rem:proper-comparison-not-circular}
The functor \(f_*\) and both canonical transformations were constructed
before any use of ordinary proper functoriality.  Classical comparison is
used only to detect their invertibility.  In particular, the exceptional
functor \(f_!\) of the next chapter will not be defined by transport from the
ordinary theory; it will be constructed from spectral compactifications and
the proper support property proved below.
\end{remark}





\begin{corollary}[Proper projection with virtual twists]
\label{cor:proper-virtual-projection}
Let \(f:X\to S\) be proper and let \(\xi\in K(S)\).  Then
\[
    f_*\Sigma^{f^*\xi}A
    \simeq
    \Sigma^\xi f_*A.
\]
\end{corollary}

\begin{proof}
Apply Theorem~\ref{thm:proper-projection-formula} to the invertible object
\(\operatorname{th}_S(\xi)\).
\end{proof}

\section{Open--proper exchange and the support property}
\label{sec:proper-excision}

\subsection{Open--proper exchange}

Consider a Cartesian square
\[
\begin{tikzcd}[column sep=large,row sep=large]
    U' \arrow[r,"j'"] \arrow[d,"q"']
        & X' \arrow[d,"p"] \\
    U \arrow[r,"j"']
        & X,
\end{tikzcd}
\]
where
\[
    p\in\mathcal P
    \qquad\text{and}\qquad
    j\in\mathcal J.
\]
Then \(j'\in\mathcal J\) and \(q\in\mathcal P\).

\begin{construction}[The open--proper exchange transformation]
\label{constr:open-proper-exchange}
Let
\[
    i:Z\hookrightarrow X
\]
be the closed complement of \(j\), and put
\[
    Z':=X'\times_XZ.
\]
Write
\[
    i':Z'\hookrightarrow X'
\]
for the pulled-back closed immersion.

For \(A\in\SH(U')\), put
\[
    E:=p_*j'_\sharp A.
\]
Proper base change for the closed Cartesian square gives
\[
\begin{aligned}
    i^*E
    &=
    i^*p_*j'_\sharp A \\
    &\simeq
    (p|_{Z'})_*i'^*j'_\sharp A \\
    &\simeq
    0
\end{aligned}
\]
by open--closed orthogonality.

The localization triangle for
\[
    Z\xrightarrow{i}X\xleftarrow{j}U
\]
is
\[
    j_\sharp j^*E
    \xrightarrow{\epsilon_j(E)}
    E
    \longrightarrow
    i_*i^*E.
\]
Since \(i^*E\simeq0\), the localization counit is an equivalence.  Its
inverse gives
\[
    E
    \xrightarrow{\simeq}
    j_\sharp j^*E.
\]

Proper base change for the original open Cartesian square gives
\[
\begin{aligned}
    j^*E
    &=
    j^*p_*j'_\sharp A \\
    &\xrightarrow{\simeq}
    q_*j'^*j'_\sharp A \\
    &\simeq
    q_*A.
\end{aligned}
\]
Define
\[
    \operatorname{Ex}^{\sharp}_{*}(A):
    p_*j'_\sharp A
    \longrightarrow
    j_\sharp q_*A
\]
to be the composite
\[
\begin{aligned}
    p_*j'_\sharp A
    &=E \\
    &\xrightarrow{\epsilon_j(E)^{-1}}
    j_\sharp j^*E \\
    &\xrightarrow{\simeq}
    j_\sharp q_*A.
\end{aligned}
\]
The construction is natural in \(A\).
\end{construction}

\begin{theorem}[Open--proper exchange]
\label{thm:open-proper-exchange}
For every Cartesian square
\[
\begin{tikzcd}[column sep=large,row sep=large]
    U' \arrow[r,"j'"] \arrow[d,"q"']
        & X' \arrow[d,"p"] \\
    U \arrow[r,"j"']
        & X
\end{tikzcd}
\]
with \(p\in\mathcal P\) and \(j\in\mathcal J\), the canonical exchange
transformation is an equivalence:
\[
    \operatorname{Ex}^{\sharp}_{*}:
    p_*j'_\sharp
    \xrightarrow{\simeq}
    j_\sharp q_*.
\]
It is compatible with arbitrary base change and with composition of proper
morphisms.
\end{theorem}

\begin{proof}
The construction of
\[
    \operatorname{Ex}^{\sharp}_{*}(A)
\]
is a composite of equivalences: proper base change and open--closed
orthogonality show that \(i^*p_*j'_\sharp A\simeq0\); localization then
identifies \(p_*j'_\sharp A\) with its extension from \(U\); and proper base
change for the open square identifies the restriction with \(q_*A\).
Hence the exchange transformation is an equivalence.

For composition, suppose that
\[
    X''\xrightarrow{p'}X'\xrightarrow{p}X
\]
are proper and that compatible Cartesian open subschemes give morphisms
\[
    U''\xrightarrow{q'}U'\xrightarrow{q}U.
\]
Both the exchange map for \(pp'\) and the composite of the exchange maps for
\(p'\) and \(p\) are transformations between objects in the essential image
of
\[
    j_\sharp:\SH(U)\longrightarrow\SH(X).
\]
After applying \(j^*\), both become the canonical direct-image composition
equivalence
\[
    (qq')_*
    \simeq
    q_*q'_*.
\]
The restriction functor
\[
    j^*:\operatorname{Im}(j_\sharp)\longrightarrow\SH(U)
\]
is an equivalence.  Therefore the two transformations agree.  The same
uniqueness argument supplies all higher composition coherences.

Every step in
Construction~\ref{constr:open-proper-exchange} is preserved by arbitrary
base change: proper base change, pullback of the closed complement,
localization, open--closed orthogonality, and the direct-image composition
constraints.  Hence the exchange equivalence is compatible with arbitrary
base change.
\end{proof}

\subsection{Proper support}

\begin{definition}[Proper excision square]
\label{def:proper-excision-square}
A proper excision square is a Cartesian square
\[
\begin{tikzcd}[column sep=large,row sep=large]
    Z' \arrow[r,"i'"] \arrow[d,"q"']
        & X' \arrow[d,"p"] \\
    Z \arrow[r,"i"']
        & X,
\end{tikzcd}
\]
where \(i\) is a closed immersion, \(p\in\mathcal P\), and, writing
\[
    j:U\hookrightarrow X,
    \qquad
    j':U'\hookrightarrow X'
\]
for the open complements, the induced map
\[
    u:U'\xrightarrow{\simeq}U
\]
is an equivalence.
\end{definition}

\begin{corollary}[Proper support]
\label{thm:proper-support}
For every proper excision square, there is a canonical equivalence
\[
    p_*j'_\sharp
    \xrightarrow{\simeq}
    j_\sharp u_*.
\]
After identifying \(u\) with the identity of \(U\), this becomes
\[
    p_*j'_\sharp
    \simeq
    j_\sharp.
\]
\end{corollary}

\begin{proof}
Apply Theorem~\ref{thm:open-proper-exchange} to the Cartesian square of open
complements.  Since \(u\) is an equivalence, so is
\[
    u_*:\SH(U')\xrightarrow{\simeq}\SH(U).
\]
\end{proof}

\begin{corollary}[Proper invariance of extension by zero]
\label{cor:proper-invariance-extension-zero}
Let
\[
    r:\overline X'\longrightarrow\overline X
\]
be proper, and suppose that there is an open spectral scheme \(X\) with
Cartesian open immersions
\[
\begin{tikzcd}[column sep=large,row sep=large]
    X \arrow[r,"j'"] \arrow[d,"\id_X"']
        & \overline X' \arrow[d,"r"] \\
    X \arrow[r,"j"']
        & \overline X.
\end{tikzcd}
\]
Then
\[
    r_*j'_\sharp
    \simeq
    j_\sharp.
\]
\end{corollary}

\begin{proof}
Take the closed complements and apply
Corollary~\ref{thm:proper-support}.
\end{proof}

\begin{proposition}[Composition coherence of support]
\label{prop:proper-support-composition}
Let
\[
    \overline X''
    \xrightarrow{r'}
    \overline X'
    \xrightarrow{r}
    \overline X
\]
be proper morphisms restricting to the identity on a common open spectral
subscheme \(X\).  Under the canonical equivalence
\[
    (rr')_*
    \simeq
    r_*r'_*,
\]
the support equivalence for \(rr'\) agrees with the composite
\[
    (rr')_*j''_\sharp
    \simeq
    r_*r'_*j''_\sharp
    \simeq
    r_*j'_\sharp
    \simeq
    j_\sharp.
\]
The same compatibility holds coherently for arbitrary finite composites.
\end{proposition}

\begin{proof}
This is the composition compatibility of open--proper exchange from
Theorem~\ref{thm:open-proper-exchange}, specialized to the common open
\(X\).  Equivalently, both transformations are maps between objects
supported on \(X\), and after applying \(j^*\) both become the identity.
Since \(j^*\) is an equivalence on the supported subcategory, they agree.
The same uniqueness gives all higher coherences.
\end{proof}

\begin{proposition}[Base change of support]
\label{prop:proper-support-base-change}
The proper support equivalence is compatible with arbitrary base change of
proper excision squares.
\end{proposition}

\begin{proof}
This is the base-change compatibility of open--proper exchange from
Theorem~\ref{thm:open-proper-exchange}, specialized to the open-complement
square.
\end{proof}

\section{Compactifiable morphisms}
\label{sec:compactifiable-morphisms}

\begin{definition}[Compactification]
\label{def:spectral-compactification}
Let
\[
    f:X\longrightarrow S
\]
belong to \(\mathcal E\).

A compactification of \(f\) is a factorization
\[
    X
    \xrightarrow{j}
    \overline X
    \xrightarrow{p}
    S
\]
together with a specified equivalence
\[
    pj\simeq f,
\]
where
\[
    j\in\mathcal J
    \qquad\text{and}\qquad
    p\in\mathcal P.
\]
The morphism \(f\) is compactifiable if it admits a compactification.
\end{definition}

\begin{definition}[Refinement of compactifications]
\label{def:compactification-refinement}
Let
\[
    c_a
    =
    \bigl(
      X\xrightarrow{j_a}\overline X_a\xrightarrow{p_a}S
    \bigr),
    \qquad
    a=1,2,
\]
be compactifications of \(f\).

A refinement
\[
    c_2\longrightarrow c_1
\]
is a morphism
\[
    r:\overline X_2\longrightarrow\overline X_1
\]
over \(S\), belonging to \(\mathcal P\), together with a specified
Cartesian square
\[
\begin{tikzcd}[column sep=large,row sep=large]
    X \arrow[r,"j_2"] \arrow[d,"\id_X"']
        & \overline X_2 \arrow[d,"r"] \\
    X \arrow[r,"j_1"']
        & \overline X_1.
\end{tikzcd}
\]
The square is required to be compatible with the specified equivalences
\[
    p_aj_a\simeq f.
\]
\end{definition}

\begin{definition}[The category of compactifications]
\label{def:category-compactifications}
Write
\[
    \operatorname{Comp}(f)
\]
for the \(\infty\)-category defined as follows.

Its objects are compactifications of \(f\).  For two objects
\[
    c_a
    =
    \bigl(
      X\xrightarrow{j_a}\overline X_a\xrightarrow{p_a}S
    \bigr),
    \qquad
    a=1,2,
\]
the mapping space
\[
    \Map_{\operatorname{Comp}(f)}(c_2,c_1)
\]
is the union of those connected components of the homotopy fibre
\[
\Map_{\operatorname{Sch}^{\mathrm{sp}}_{/S}}
    (\overline X_2,\overline X_1)
\times_{
\Map_{\operatorname{Sch}^{\mathrm{sp}}_{/S}}
    (X,\overline X_1)
}
\{j_1\}
\]
whose underlying morphism
\[
    r:\overline X_2\longrightarrow\overline X_1
\]
belongs to \(\mathcal P\) and for which the induced square
\[
\begin{tikzcd}[column sep=large,row sep=large]
    X \arrow[r,"j_2"] \arrow[d,"\id_X"']
        & \overline X_2 \arrow[d,"r"] \\
    X \arrow[r,"j_1"']
        & \overline X_1
\end{tikzcd}
\]
is Cartesian.  The map to
\[
\Map_{\operatorname{Sch}^{\mathrm{sp}}_{/S}}
    (X,\overline X_1)
\]
is induced by precomposition with \(j_2\).

Composition and all higher simplices are inherited from the slice
\(\infty\)-category
\[
    \operatorname{Sch}^{\mathrm{sp}}_{/S}.
\]
\end{definition}

\begin{proposition}[Elementary permanence of compactifiability]
\label{prop:compactifiable-permanence}
Compactifiable morphisms in \(\mathcal E\) are stable under the following
operations:

\begin{enumerate}
    \item arbitrary base change;

    \item base restriction along a quasi-compact open subscheme of the
    target;

    \item precomposition with a quasi-compact open immersion;

    \item products over a common base;

    \item postcomposition with a morphism in
    \(\mathcal P\cap\mathcal E\).
\end{enumerate}
\end{proposition}

\begin{proof}
Let
\[
    X
    \xrightarrow{j}
    \overline X
    \xrightarrow{p}
    S
\]
be a compactification.

Arbitrary base change preserves quasi-compact open immersions and carries
morphisms in \(\mathcal P\) to morphisms in \(\mathcal P\).  Pulling back
the displayed factorization therefore gives a compactification after
arbitrary base change.  Base restriction to a quasi-compact open subscheme
of \(S\) is a special case.

If
\[
    v:V\hookrightarrow X
\]
is a quasi-compact open immersion, then
\[
    V
    \xrightarrow{jv}
    \overline X
    \xrightarrow{p}
    S
\]
is a compactification of \(fv\), because a composite of quasi-compact open
immersions is quasi-compact open.

Let
\[
    f_a:X_a\longrightarrow S,
    \qquad
    a=1,2,
\]
be compactifiable, with compactifications
\[
    X_a
    \hookrightarrow
    \overline X_a
    \longrightarrow
    S.
\]
Then
\[
    X_1\times_SX_2
    \hookrightarrow
    \overline X_1\times_S\overline X_2
    \longrightarrow
    S
\]
is a compactification of
\[
    X_1\times_SX_2\longrightarrow S.
\]
The first map is quasi-compact open and the second belongs to
\(\mathcal P\).

Suppose
\[
    r:S\longrightarrow T
\]
belongs to \(\mathcal P\cap\mathcal E\).  Then
\[
    X
    \hookrightarrow
    \overline X
    \xrightarrow{p}
    S
    \xrightarrow{r}
    T
\]
is a compactification of \(rf\), because \(rp\) belongs to
\(\mathcal P\).
\end{proof}

\section{Spectral Nagata compactification}
\label{sec:spectral-nagata}

\begin{proposition}[Spectral cobase change along a nil-immersion]
\label{prop:spectral-nil-cobase-change-open}
Consider a diagram of spectral schemes
\[
\begin{tikzcd}[column sep=large,row sep=large]
    X_0 \arrow[r,"i"] \arrow[d,"j_0"']
        & X \\
    Y_0,
\end{tikzcd}
\]
where \(X_0\) and \(Y_0\) are ordinary schemes regarded as spectral
schemes, and

\begin{enumerate}
    \item \(i:X_0\hookrightarrow X\) is a closed immersion inducing an
    equivalence on classical truncations;

    \item \(j_0:X_0\hookrightarrow Y_0\) is a quasi-compact open immersion.
\end{enumerate}

Then the pushout
\[
    Y
    :=
    X\amalg_{X_0}Y_0
\]
exists in spectral schemes, and the resulting square
\[
\begin{tikzcd}[column sep=large,row sep=large]
    X_0 \arrow[r,"i"] \arrow[d,"j_0"']
        & X \arrow[d,"j"] \\
    Y_0 \arrow[r,"i_Y"']
        & Y
\end{tikzcd}
\]
is Cartesian as well as cocartesian.  Moreover,

\begin{enumerate}
    \item \(j:X\hookrightarrow Y\) is a quasi-compact open immersion;

    \item \(i_Y:Y_0\hookrightarrow Y\) is a closed immersion inducing an
    equivalence
    \[
        (Y_0)_{\mathrm{cl}}
        \xrightarrow{\simeq}
        Y_{\mathrm{cl}};
    \]

    \item the construction is compatible with arbitrary base change;

    \item if \(X\) and \(Y_0\) are quasi-compact and quasi-separated, then
    \(Y\) is quasi-compact and quasi-separated.
\end{enumerate}
\end{proposition}

\begin{proof}
We use the spectral nil-cobase-change theorem in its universally
base-change-stable form; see
\cite{SpectralNilCobaseChange}.  In the present situation it asserts that
the pushout
\[
    Y
    :=
    X\amalg_{X_0}Y_0
\]
exists in spectral schemes, that the induced morphism
\[
    j:X\longrightarrow Y
\]
is the open immersion obtained by extending the classical open
\[
    j_0:X_0\hookrightarrow Y_0,
\]
and that the construction is preserved by arbitrary base change.  It also
identifies the structure sheaf, on the Zariski site of \(Y_0\), by the
pullback formula
\[
    \mathcal O_Y
    \simeq
    \mathcal O_{Y_0}
    \times_{j_{0*}\mathcal O_{X_0}}
    j_{0*}\mathcal O_X.
\]

Because
\[
    i:X_0\hookrightarrow X
\]
induces an equivalence on classical truncations, the underlying topological
space of \(X_0\) is canonically identified with that of \(X\).  We therefore
view this space as the open subset
\[
    |X_0|\subseteq|Y_0|.
\]
Pulling the displayed structure-sheaf formula back along \(j_0\) gives
\[
\begin{aligned}
    j_0^*\mathcal O_Y
    &\simeq
    \mathcal O_{X_0}
    \times_{\mathcal O_{X_0}}
    \mathcal O_X \\
    &\simeq
    \mathcal O_X.
\end{aligned}
\]
Thus
\[
    j:X\longrightarrow Y
\]
is the open immersion corresponding to
\[
    |X_0|\subseteq|Y_0|.
\]
It is quasi-compact because \(j_0\) is quasi-compact.

Let \(F\) be the fibre of
\[
    \mathcal O_X
    \longrightarrow
    \mathcal O_{X_0}.
\]
Since \(i\) induces an equivalence on classical truncations, the induced map
on \(\pi_0\) is an equivalence, and therefore
\[
    \pi_0(F)=0.
\]
Taking fibres in the structure-sheaf pullback square identifies the fibre of
\[
    \mathcal O_Y
    \longrightarrow
    \mathcal O_{Y_0}
\]
with
\[
    j_{0*}F.
\]
Consequently,
\[
    \pi_0(\mathcal O_Y)
    \simeq
    \pi_0(\mathcal O_{Y_0}),
\]
and hence
\[
    Y_{\mathrm{cl}}
    \simeq
    (Y_0)_{\mathrm{cl}}.
\]
The induced morphism
\[
    i_Y:Y_0\longrightarrow Y
\]
is therefore a closed nil-immersion inducing an equivalence
\[
    (Y_0)_{\mathrm{cl}}
    \xrightarrow{\simeq}
    Y_{\mathrm{cl}}.
\]

It remains to identify the pullback of the open immersion \(j\) along
\(i_Y\).  The fibre product
\[
    Y_0\times_YX
\]
is the restriction of \(Y_0\) to the open spectral subscheme
\[
    X\subseteq Y.
\]
Its underlying classical open is \(X_0\), and the restriction of its
structure sheaf is
\[
    \mathcal O_{X_0}.
\]
Therefore
\[
    Y_0\times_YX
    \simeq
    X_0.
\]
Thus the defining square is Cartesian.  It is cocartesian by the definition
of \(Y\).

Compatibility with arbitrary base change is part of the universally
base-change-stable nil-cobase-change theorem invoked in the first paragraph.
In particular, after any morphism
\[
    T\longrightarrow Y,
\]
the pulled-back square is canonically the corresponding nil-cobase pushout.

Finally, quasi-compactness and quasi-separatedness of spectral schemes are
detected on classical truncations.  Since
\[
    Y_{\mathrm{cl}}
    \simeq
    (Y_0)_{\mathrm{cl}},
\]
qcqs-ness of \(Y_0\) implies qcqs-ness of \(Y\).
\end{proof}

\begin{theorem}[Spectral Nagata compactification]
\label{thm:spectral-nagata-compactification}
Let
\[
    f:X\longrightarrow S
\]
be a separated morphism of finite presentation between quasi-compact
quasi-separated spectral schemes.  Then there exists a factorization
\[
    X
    \xrightarrow{j}
    \overline X
    \xrightarrow{p}
    S
\]
such that
\[
    j\in\mathcal J
    \qquad\text{and}\qquad
    p\in\mathcal P.
\]
Thus \(j\) is a quasi-compact open immersion and
\[
    p_{\mathrm{cl}}:
    \overline X_{\mathrm{cl}}
    \longrightarrow
    S_{\mathrm{cl}}
\]
is proper.  The compactifying object \(\overline X\) is a spectral scheme.
\end{theorem}

\begin{proof}
Classical truncation preserves finite presentation and separatedness.
Consequently,
\[
    f_{\mathrm{cl}}:
    X_{\mathrm{cl}}
    \longrightarrow
    S_{\mathrm{cl}}
\]
is a separated morphism of finite presentation between quasi-compact
quasi-separated ordinary schemes.

Classical Nagata compactification gives a factorization
\[
    X_{\mathrm{cl}}
    \xrightarrow{j_0}
    \overline X_0
    \xrightarrow{p_0}
    S_{\mathrm{cl}},
\]
where \(j_0\) is a quasi-compact open immersion and \(p_0\) is proper.

Regard \(\overline X_0\) as a discrete spectral scheme and regard \(p_0\)
as a morphism
\[
    \overline X_0
    \longrightarrow
    S_{\mathrm{cl}}
    \longrightarrow
    S.
\]
The classical truncation morphism
\[
    \iota_X:
    X_{\mathrm{cl}}
    \hookrightarrow
    X
\]
is a closed immersion inducing an equivalence on classical truncations.
Apply
Proposition~\ref{prop:spectral-nil-cobase-change-open}
to the diagram
\[
\begin{tikzcd}[column sep=large,row sep=large]
    X_{\mathrm{cl}}
        \arrow[r,"\iota_X"]
        \arrow[d,"j_0"']
        & X \\
    \overline X_0.
\end{tikzcd}
\]
It produces a spectral scheme
\[
    \overline X
    :=
    X\amalg_{X_{\mathrm{cl}}}\overline X_0
\]
and a Cartesian and cocartesian square
\[
\begin{tikzcd}[column sep=large,row sep=large]
    X_{\mathrm{cl}}
        \arrow[r,"\iota_X"]
        \arrow[d,"j_0"']
        & X \arrow[d,"j"] \\
    \overline X_0
        \arrow[r]
        & \overline X.
\end{tikzcd}
\]
The morphism
\[
    j:X\hookrightarrow\overline X
\]
is a quasi-compact open immersion.

The compatible maps
\[
    X\longrightarrow S
    \qquad\text{and}\qquad
    \overline X_0\longrightarrow S
\]
induce a morphism
\[
    p:\overline X\longrightarrow S.
\]
The nil-cobase-change theorem gives a canonical equivalence
\[
    \overline X_{\mathrm{cl}}
    \simeq
    \overline X_0
\]
under which
\[
    p_{\mathrm{cl}}
    \simeq
    p_0.
\]
Since \(p_0\) is proper, \(p\) belongs to \(\mathcal P\).  Moreover,
\(\overline X\) is quasi-compact and quasi-separated by
Proposition~\ref{prop:spectral-nil-cobase-change-open}.

Thus
\[
    X
    \xrightarrow{j}
    \overline X
    \xrightarrow{p}
    S
\]
is the required spectral compactification.
\end{proof}

\begin{corollary}[Every exceptional morphism is compactifiable]
\label{cor:every-E-compactifiable}
Every morphism in \(\mathcal E\) is compactifiable.
\end{corollary}

\begin{proof}
This is Theorem~\ref{thm:spectral-nagata-compactification}.
\end{proof}

\begin{corollary}[Composition of compactifiable morphisms]
\label{cor:compactifiable-composition}
Let
\[
    X\xrightarrow{f}Y\xrightarrow{g}S
\]
be morphisms in \(\mathcal E\).  Then the composite
\[
    gf:X\longrightarrow S
\]
is compactifiable.
\end{corollary}

\begin{proof}
The composite \(gf\) belongs to \(\mathcal E\), because separated morphisms
of finite presentation are stable under composition.  Apply
Theorem~\ref{thm:spectral-nagata-compactification}.
\end{proof}

\subsection{Classical strict refinements and spectral lifting}

Put
\[
    X_0:=X_{\mathrm{cl}},
    \qquad
    S_0:=S_{\mathrm{cl}},
    \qquad
    f_0:=f_{\mathrm{cl}}.
\]

\begin{definition}[The strict classical compactification category]
\label{def:strict-classical-compactifications}
Let
\[
    \operatorname{Comp}^{\square}_{\mathrm{cl}}(f_0)
\]
be the ordinary category, regarded through its nerve as an
\(\infty\)-category, whose objects are classical compactifications
\[
    X_0
    \xrightarrow{j_0}
    \overline X_0
    \xrightarrow{p_0}
    S_0
\]
and whose morphisms
\[
    c_2\longrightarrow c_1
\]
are proper morphisms
\[
    r_0:\overline X_{0,2}\longrightarrow\overline X_{0,1}
\]
over \(S_0\) for which the square
\[
\begin{tikzcd}[column sep=large,row sep=large]
    X_0 \arrow[r,"j_{0,2}"] \arrow[d,"\id_{X_0}"']
        & \overline X_{0,2} \arrow[d,"r_0"] \\
    X_0 \arrow[r,"j_{0,1}"']
        & \overline X_{0,1}
\end{tikzcd}
\]
is Cartesian.
\end{definition}

\begin{definition}[Dense compactification]
\label{def:dense-classical-compactification}
A classical compactification
\[
    X_0\hookrightarrow\overline X_0\longrightarrow S_0
\]
is dense if \(X_0\) is schematically dense in \(\overline X_0\).
\end{definition}

\begin{lemma}[Dense refinements and uniqueness]
\label{lem:dense-compactification-uniqueness}
Every object of
\(\operatorname{Comp}^{\square}_{\mathrm{cl}}(f_0)\) admits a dense
refinement.  If \(\overline X_2\) is a dense compactification, then any two
refinements
\[
    r,s:\overline X_2\rightrightarrows\overline X_1
\]
are equal.
\end{lemma}

\begin{proof}
Let
\[
    X_0\hookrightarrow\overline X_0\longrightarrow S_0
\]
be a compactification.  Let \(\overline X_0^{\mathrm{dens}}\) be the
scheme-theoretic closure of \(X_0\) in \(\overline X_0\).  Since \(X_0\)
is open in \(\overline X_0\), it remains open in its closure.  The closure
is closed in a proper \(S_0\)-scheme and is therefore proper over \(S_0\).
Its inverse image over the open \(X_0\subseteq\overline X_0\) is exactly
\(X_0\).  Hence
\[
    \overline X_0^{\mathrm{dens}}
    \longrightarrow
    \overline X_0
\]
is a dense refinement.

Now let \(r,s:\overline X_2\rightrightarrows\overline X_1\) be refinements
with \(X_0\) schematically dense in \(\overline X_2\).  Since
\(\overline X_1\to S_0\) is separated, the equalizer of \(r\) and \(s\)
is a closed subscheme of \(\overline X_2\).  It contains \(X_0\), because
both maps restrict to the identity there.  Schematic density forces the
equalizer to be all of \(\overline X_2\).  Thus \(r=s\).
\end{proof}

\begin{proposition}[Cofilteredness of strict classical compactifications]
\label{prop:classical-strict-compactifications-cofiltered}
The category
\[
    \operatorname{Comp}^{\square}_{\mathrm{cl}}(f_0)
\]
is nonempty and cofiltered.
\end{proposition}

\begin{proof}
Nonemptiness follows from classical Nagata compactification.

Consider finitely many dense compactifications
\[
    X_0\hookrightarrow\overline X_a\longrightarrow S_0,
    \qquad
    1\leq a\leq n.
\]
Form the proper \(S_0\)-scheme
\[
    P
    :=
    \overline X_1
    \times_{S_0}
    \cdots
    \times_{S_0}
    \overline X_n.
\]
The open subscheme
\[
    U
    :=
    X_0
    \times_{S_0}
    \cdots
    \times_{S_0}
    X_0
\]
is open in \(P\).  The multi-diagonal
\[
    \delta:X_0\longrightarrow U
\]
is a closed immersion because \(f_0\) is separated.  Hence
\(\delta:X_0\to P\) is a locally closed immersion.

Let \(\overline X\) be the scheme-theoretic closure of \(\delta(X_0)\) in
\(P\).  Then
\[
    X_0\hookrightarrow\overline X
\]
is open and schematically dense, and \(\overline X\to S_0\) is proper.
For each \(a\), let
\[
    \rho_a:\overline X\longrightarrow\overline X_a
\]
be the corresponding projection.  It is proper.

Put
\[
    W_a
    :=
    \overline X\times_{\overline X_a}X_0.
\]
The morphism \(W_a\to X_0\) is proper, and the diagonal copy of \(X_0\)
inside \(W_a\) is an open section.  Since \(W_a\to X_0\) is separated,
that section is also closed.  Schematic density of \(X_0\) in
\(\overline X\) restricts to schematic density in the open subscheme
\(W_a\).  Therefore
\[
    W_a\simeq X_0.
\]
Thus every \(\rho_a\) is a strict refinement, so every finite family has a
common dense refinement.

Let now
\[
    D:K\longrightarrow
    \operatorname{Comp}^{\square}_{\mathrm{cl}}(f_0)
\]
be a diagram indexed by a finite simplicial set.  Choose dense refinements of
its finitely many vertices and then a common dense refinement
\(\overline X\).  For every arrow of \(D\), the two induced maps from
\(\overline X\) to its target agree on \(X_0\), and hence agree by
Lemma~\ref{lem:dense-compactification-uniqueness}.  All higher compatibility
conditions follow from the same uniqueness.  Hence \(\overline X\) defines a
cone over \(D\), proving cofilteredness.
\end{proof}

\begin{construction}[Truncation and nil-cobase lifting]
\label{constr:compactification-lifting-adjunction}
Classical truncation defines a functor
\[
    \tau_f:
    \operatorname{Comp}(f)
    \longrightarrow
    \operatorname{Comp}^{\square}_{\mathrm{cl}}(f_0).
\]

For a classical compactification
\[
    c_0=
    \bigl(
      X_0\xrightarrow{j_0}\overline X_0\xrightarrow{p_0}S_0
    \bigr),
\]
define
\[
    \lambda_f(c_0)
    :=
    \left(
      X
      \longrightarrow
      X\amalg_{X_0}\overline X_0
      \longrightarrow
      S
    \right).
\]
Proposition~\ref{prop:spectral-nil-cobase-change-open} shows that the first
map is a quasi-compact open immersion, that the middle spectral scheme is
qcqs, and that its classical truncation is \(\overline X_0\).  The second
map therefore belongs to \(\mathcal P\).

A strict classical refinement
\[
    r_0:\overline X_{0,2}\longrightarrow\overline X_{0,1}
\]
induces, by the pushout universal property, a morphism
\[
    \lambda_f(r_0):
    X\amalg_{X_0}\overline X_{0,2}
    \longrightarrow
    X\amalg_{X_0}\overline X_{0,1}.
\]
Its classical truncation is \(r_0\), so it belongs to \(\mathcal P\).  The
pullback of the target open is an open spectral subscheme whose classical
open is \(X_0\); uniqueness of open spectral subschemes with a prescribed
classical open identifies it with \(X\).  Thus \(\lambda_f(r_0)\) is a
refinement.  This defines a functor
\[
    \lambda_f:
    \operatorname{Comp}^{\square}_{\mathrm{cl}}(f_0)
    \longrightarrow
    \operatorname{Comp}(f).
\]
\end{construction}

\begin{proposition}[The compactification lifting adjunction]
\label{prop:compactification-lifting-adjunction}
There is an adjunction
\[
    \lambda_f:
    \operatorname{Comp}^{\square}_{\mathrm{cl}}(f_0)
    \rightleftarrows
    \operatorname{Comp}(f):
    \tau_f.
\]
Its unit is an equivalence:
\[
    \tau_f\lambda_f
    \simeq
    \id.
\]
For a spectral compactification
\[
    c=
    \bigl(
      X\xrightarrow{j}\overline X\xrightarrow{p}S
    \bigr),
\]
the counit is the canonical refinement
\[
    \epsilon_c:
    X\amalg_{X_0}\overline X_{\mathrm{cl}}
    \longrightarrow
    \overline X
\]
induced by \(j\) and the truncation immersion
\(\overline X_{\mathrm{cl}}\to\overline X\).
\end{proposition}

\begin{proof}
Let
\[
    c_0=
    \bigl(
      X_0\hookrightarrow\overline X_0\to S_0
    \bigr)
\]
be classical and let
\[
    c=
    \bigl(
      X\hookrightarrow\overline X\to S
    \bigr)
\]
be spectral.  The pushout universal property gives
\[
\begin{aligned}
    \Map_{/S}
    \bigl(
      X\amalg_{X_0}\overline X_0,
      \overline X
    \bigr)
    \simeq
    &\Map_{/S}(X,\overline X) \\
    &\times_{\Map_{/S}(X_0,\overline X)}
    \Map_{/S}(\overline X_0,\overline X).
\end{aligned}
\]
Taking the homotopy fibre over the specified map
\(j:X\to\overline X\) fixes the first factor.

For every ordinary scheme \(V_0\) and spectral scheme \(W\), the adjunction
between discrete enhancement and classical truncation gives
\[
    \Map_{\operatorname{Sch}^{\mathrm{sp}}}(V_0,W)
    \simeq
    \Map_{\operatorname{Sch}}(V_0,W_{\mathrm{cl}}).
\]
The same equivalence holds in the relevant slice categories.  It follows
that the preceding homotopy fibre is canonically the discrete mapping space
of maps
\[
    \overline X_0
    \longrightarrow
    \overline X_{\mathrm{cl}}
\]
over \(S_0\) restricting to the identity on \(X_0\).

Properness is detected on classical truncations.  The Cartesian open-square
condition is also detected classically: if
\[
    V
    :=
    \bigl(
      X\amalg_{X_0}\overline X_0
    \bigr)
    \times_{\overline X}X,
\]
then \(V\) is an open spectral subscheme and
\[
    V_{\mathrm{cl}}
    \simeq
    \overline X_0
    \times_{\overline X_{\mathrm{cl}}}
    X_0.
\]
If the classical square is Cartesian, this classical open is \(X_0\), so
\(V\simeq X\).  Therefore
\[
    \Map_{\operatorname{Comp}(f)}
    \bigl(
      \lambda_f(c_0),c
    \bigr)
    \simeq
    \Map_{\operatorname{Comp}^{\square}_{\mathrm{cl}}(f_0)}
    \bigl(
      c_0,\tau_f(c)
    \bigr).
\]
This equivalence is natural in both variables and proves the adjunction.

Proposition~\ref{prop:spectral-nil-cobase-change-open} gives
\[
    \tau_f\lambda_f(c_0)
    \simeq
    c_0,
\]
so the unit is an equivalence.  The counit is induced by the compatible maps
\[
    X\longrightarrow\overline X
    \qquad\text{and}\qquad
    \overline X_{\mathrm{cl}}
    \longrightarrow
    \overline X.
\]
Its classical truncation is the identity of
\(\overline X_{\mathrm{cl}}\), so it belongs to \(\mathcal P\).  Its open
square is Cartesian by the classical-open detection argument above.  Hence
it is a refinement in \(\operatorname{Comp}(f)\).
\end{proof}

\begin{proposition}[Refinement of a finite compactification diagram]
\label{prop:finite-compactification-diagram-refinement}
Let
\[
    D:K\longrightarrow\operatorname{Comp}(f)
\]
be a diagram indexed by a finite simplicial set.  Then \(D\) admits a
coherently compatible refining cone.
\end{proposition}

\begin{proof}
Apply classical truncation to obtain
\[
    \tau_fD:
    K\longrightarrow
    \operatorname{Comp}^{\square}_{\mathrm{cl}}(f_0).
\]
By Proposition~\ref{prop:classical-strict-compactifications-cofiltered},
there is a cone
\[
    c_0\longrightarrow\tau_fD.
\]
Applying \(\lambda_f\) gives a cone
\[
    \lambda_f(c_0)
    \longrightarrow
    \lambda_f\tau_fD.
\]
Compose objectwise with the counit of
Proposition~\ref{prop:compactification-lifting-adjunction}:
\[
    \lambda_f\tau_fD
    \longrightarrow
    D.
\]
Naturality of the counit supplies all edge and higher-simplex
compatibilities.  The resulting composite is the required refining cone.
\end{proof}

\begin{theorem}[Contractibility of compactifications]
\label{thm:compactification-category-contractible}
For every morphism
\[
    f:X\longrightarrow S
\]
in \(\mathcal E\), the \(\infty\)-category
\[
    \operatorname{Comp}(f)
\]
is nonempty and cofiltered.  Consequently, its groupoid completion is
contractible:
\[
    \bigl|
      \operatorname{Comp}(f)
    \bigr|
    \simeq
    *.
\]
\end{theorem}

\begin{proof}
Nonemptiness follows from
Theorem~\ref{thm:spectral-nagata-compactification}.  Cofilteredness is
Proposition~\ref{prop:finite-compactification-diagram-refinement}.  The
groupoid completion of a nonempty cofiltered \(\infty\)-category is
contractible.
\end{proof}

\begin{proposition}[The compactification diagram of direct images]
\label{prop:compactification-direct-image-diagram}
Let
\[
    f:X\longrightarrow S
\]
belong to \(\mathcal E\).

The assignment
\[
    \bigl(
      X\xrightarrow{j}\overline X\xrightarrow{p}S
    \bigr)
    \longmapsto
    p_*j_\sharp
\]
extends to a functor
\[
    \Phi_f:
    \operatorname{Comp}(f)
    \longrightarrow
    \Fun^L
    \bigl(
      \SH(X),
      \SH(S)
    \bigr).
\]

For every refinement
\[
    r:
    \bigl(
      X\xrightarrow{j_2}\overline X_2\xrightarrow{p_2}S
    \bigr)
    \longrightarrow
    \bigl(
      X\xrightarrow{j_1}\overline X_1\xrightarrow{p_1}S
    \bigr),
\]
the corresponding transformation
\[
    \Phi_f(r):
    p_{2*}j_{2\sharp}
    \longrightarrow
    p_{1*}j_{1\sharp}
\]
is an equivalence.

Consequently, \(\Phi_f\) factors through the contractible groupoid
completion
\[
    \bigl|
      \operatorname{Comp}(f)
    \bigr|
    \simeq
    *.
\]
For every compactification \(c\), the canonical cocone map is an
equivalence
\[
    \Phi_f(c)
    \xrightarrow{\simeq}
    \operatorname*{colim}_{d\in\operatorname{Comp}(f)}
    \Phi_f(d).
\]
Consequently, its inverse determines an equivalence
\[
    \operatorname*{colim}_{d\in\operatorname{Comp}(f)}
    \Phi_f(d)
    \xrightarrow{\simeq}
    \Phi_f(c)
\]
through a contractible space of compatible inverses.
\end{proposition}

\begin{proof}
The functor \(j_\sharp\) preserves small colimits, and
Theorem~\ref{thm:arbitrary-direct-image-cocontinuous} shows that \(p_*\)
preserves small colimits.  Thus \(p_*j_\sharp\) belongs to
\(\Fun^L(\SH(X),\SH(S))\).

Let
\[
    r:\overline X_2\longrightarrow\overline X_1
\]
be a refinement.  There is a canonical equivalence
\[
    p_{2*}
    \simeq
    p_{1*}r_*.
\]
The refinement square is an open--proper Cartesian square, and
Corollary~\ref{cor:proper-invariance-extension-zero} gives
\[
    r_*j_{2\sharp}
    \xrightarrow{\simeq}
    j_{1\sharp}.
\]
Define \(\Phi_f(r)\) by the composite
\[
\begin{aligned}
    p_{2*}j_{2\sharp}
    &\simeq
    p_{1*}r_*j_{2\sharp} \\
    &\xrightarrow{\simeq}
    p_{1*}j_{1\sharp}.
\end{aligned}
\]

Composition and all higher coherences follow from
Proposition~\ref{prop:proper-support-composition}.  Hence \(\Phi_f\) is a
functor.  Every morphism is carried to an equivalence, so it factors through
the groupoid completion.  By
Theorem~\ref{thm:compactification-category-contractible}, this groupoid is
contractible.  For a diagram of equivalences indexed by a contractible groupoid, the
canonical cocone map from every value to the colimit is an equivalence.  The
space of inverses to these compatible equivalences is contractible.
\end{proof}

\begin{corollary}[Common and equalizing refinements]
\label{cor:compactification-common-refinement}
Every finite family of compactifications of \(f\) admits a common refinement.

More generally, every finite diagram of compactifications and refinements
admits a coherently compatible refining cone.  In particular, any pair of
parallel refinements becomes homotopic after passage to a further common
refinement, and the same assertion holds for every finite higher-coherence
diagram.
\end{corollary}

\begin{proof}
Apply
Proposition~\ref{prop:finite-compactification-diagram-refinement} to the
corresponding finite discrete diagram, parallel-pair diagram, or finite
simplicial coherence diagram.
\end{proof}

\begin{proposition}[Base change of compactifications]
\label{prop:compactification-base-change}
Let
\[
    g:T\longrightarrow S
\]
be arbitrary.  Base change defines a functor
\[
    g^*:
    \operatorname{Comp}(f)
    \longrightarrow
    \operatorname{Comp}(f_T),
\]
where
\[
    f_T:X\times_ST\longrightarrow T.
\]
This functor preserves refinements and common-refinement diagrams.  It is
compatible with composition in \(g\) through coherent equivalences.
\end{proposition}

\begin{proof}
Open immersions and proper morphisms are stable under base change, so the
pullback of a compactification is a compactification.  Proper refinements
remain proper and restrict to the identity on the pulled-back open.  The
coherence is that of spectral fibre products.
\end{proof}

\begin{definition}[The category of composition data]
\label{def:compactification-composition-data}
Let
\[
    X\xrightarrow{f}Y\xrightarrow{g}S
\]
belong to \(\mathcal E\).

Write
\[
    \operatorname{Comp}(f,g)
\]
for the full \(\infty\)-subcategory of the \(\infty\)-category of
commutative diagrams
\[
\begin{tikzcd}[column sep=large,row sep=large]
    X \arrow[r,hook,"j_f"] \arrow[dr,"f"']
        & \overline X \arrow[d,"p_f"] \arrow[r,hook,"k"]
        & \overline X' \arrow[d,"r"] \\
    & Y \arrow[r,hook,"j_g"']
        & \overline Y \arrow[d,"p_g"] \\
    && S
\end{tikzcd}
\]
in which
\[
    j_f,
    \ j_g,
    \ k
    \in
    \mathcal J,
\]
\[
    p_f,
    \ r,
    \ p_g
    \in
    \mathcal P,
\]
and the central square
\[
\begin{tikzcd}[column sep=large,row sep=large]
    \overline X \arrow[r,hook,"k"] \arrow[d,"p_f"']
        & \overline X' \arrow[d,"r"] \\
    Y \arrow[r,hook,"j_g"']
        & \overline Y
\end{tikzcd}
\]
is Cartesian.

Its morphisms are commutative diagrams between such objects whose maps on
\(\overline X\), \(\overline X'\), and \(\overline Y\) belong to
\(\mathcal P\), restrict to the identity on \(X\) and \(Y\), and make all
corresponding open squares Cartesian.  Mapping spaces and higher simplices
are inherited from the corresponding diagram \(\infty\)-category of
spectral schemes.
\end{definition}

\begin{definition}[Strict classical composition data]
\label{def:strict-classical-composition-data}
Put
\[
    X_0:=X_{\mathrm{cl}},
    \qquad
    Y_0:=Y_{\mathrm{cl}},
    \qquad
    S_0:=S_{\mathrm{cl}}.
\]
Let
\[
    \operatorname{Comp}^{\square}_{\mathrm{cl}}(f_0,g_0)
\]
be the ordinary category of classical diagrams satisfying the conditions of
Definition~\ref{def:compactification-composition-data}, with proper strict
refinements, regarded as an \(\infty\)-category through its nerve.
An object is dense if
\[
    X_0\subseteq\overline X_0,
    \qquad
    \overline X_0\subseteq\overline X'_0,
    \qquad
    Y_0\subseteq\overline Y_0
\]
are schematically dense open immersions.
\end{definition}

\begin{proposition}[Cofilteredness of classical composition data]
\label{prop:classical-composition-data-cofiltered}
The category
\[
    \operatorname{Comp}^{\square}_{\mathrm{cl}}(f_0,g_0)
\]
is nonempty and cofiltered.
\end{proposition}

\begin{proof}
Choose a classical compactification
\[
    X_0\hookrightarrow\overline X_0\longrightarrow Y_0
\]
of \(f_0\), and a classical compactification
\[
    Y_0\hookrightarrow\overline Y_0\longrightarrow S_0
\]
of \(g_0\).  The composite
\[
    \overline X_0\longrightarrow Y_0\hookrightarrow\overline Y_0
\]
is separated and of finite type.  Choose a dense Nagata compactification
\[
    \overline X_0
    \hookrightarrow
    \overline X'_0
    \longrightarrow
    \overline Y_0.
\]
Put
\[
    W
    :=
    \overline X'_0\times_{\overline Y_0}Y_0.
\]
The open immersion
\[
    \overline X_0\hookrightarrow W
\]
is proper: its source is proper over \(Y_0\), its target is separated over
\(Y_0\), and the map is a \(Y_0\)-morphism.  It is therefore also closed.
Since \(\overline X_0\) is schematically dense in \(\overline X'_0\), it is
schematically dense in the open subscheme \(W\).  Hence
\[
    W\simeq\overline X_0.
\]
The central square is Cartesian, proving nonemptiness.

Every composition datum admits a dense refinement.  Indeed, first
replace \(\overline Y_0\) by the scheme-theoretic closure
\(\overline Y_0^{\mathrm{dens}}\) of \(Y_0\).  Pull back
\(\overline X'_0\) to this closed subscheme.  Next replace
\(\overline X_0\) by the scheme-theoretic closure
\(\overline X_0^{\mathrm{dens}}\) of \(X_0\), and let
\(\overline X_0'{}^{\mathrm{dens}}\) be the scheme-theoretic closure of
\(\overline X_0^{\mathrm{dens}}\) in the pulled-back proper scheme.  The
pullback
\[
    \overline X_0'{}^{\mathrm{dens}}
    \times_{\overline Y_0^{\mathrm{dens}}}
    Y_0
\]
contains \(\overline X_0^{\mathrm{dens}}\) as an open immersion.  That open
immersion is proper, because its source is proper over \(Y_0\) and its
target is separated over \(Y_0\).  It is therefore closed, and schematic
density identifies the pullback with
\(\overline X_0^{\mathrm{dens}}\).  Thus the central square remains
Cartesian.

Maps out of a dense composition datum are unique whenever they exist.
Indeed, two maps on \(\overline Y_0\) agree because they agree on the dense
open \(Y_0\) and the target is separated; two maps on \(\overline X_0\)
agree because they agree on the dense open \(X_0\); and two maps on
\(\overline X'_0\) agree because they agree on the dense open
\(\overline X_0\).

It remains to construct common dense refinements.  Consider finitely many
dense composition data indexed by \(a\).  By
Proposition~\ref{prop:classical-strict-compactifications-cofiltered}, choose
a common dense refinement
\[
    Y_0\hookrightarrow\overline Y\longrightarrow S_0
\]
of the compactifications \(Y_0\hookrightarrow\overline Y_a\).  Likewise,
choose a common dense refinement
\[
    X_0\hookrightarrow\overline X\longrightarrow Y_0
\]
of the compactifications \(X_0\hookrightarrow\overline X_a\).

Put
\[
    Q_a
    :=
    \overline X'_a
    \times_{\overline Y_a}
    \overline Y.
\]
Each \(Q_a\to\overline Y\) is proper, and its pullback to \(Y_0\) is
\(\overline X_a\).  The refinement maps
\(\overline X\to\overline X_a\) therefore induce a morphism
\[
    \overline X
    \longrightarrow
    Q
    :=
    \prod_{\overline Y}Q_a.
\]
Choose a dense compactification
\[
    \overline X
    \hookrightarrow
    B
    \longrightarrow
    \overline Y
\]
of the separated finite-type morphism
\(\overline X\to\overline Y\).  The graph of the morphism
\(\overline X\to Q\) gives a locally closed immersion
\[
    \overline X
    \longrightarrow
    B\times_{\overline Y}Q.
\]
Let \(\overline X'\) be its scheme-theoretic closure.  Then
\[
    \overline X\hookrightarrow\overline X'
\]
is open and schematically dense, and
\(\overline X'\to\overline Y\) is proper.  The projections give proper maps
from \(\overline X'\) to every \(\overline X'_a\).  Their inverse images of
\(\overline X_a\) are \(\overline X\): indeed,
\[
\begin{aligned}
    \overline X'
    \times_{\overline X'_a}
    \overline X_a
    &\simeq
    \overline X'
    \times_{\overline Y_a}
    Y_0 \\
    &\simeq
    \overline X'
    \times_{\overline Y}
    Y_0.
\end{aligned}
\]
The final term is identified with \(\overline X\) below.  Hence these maps,
together with the strict refinements on \(\overline X\) and
\(\overline Y\), define refinements of composition data.

Finally, put
\[
    W
    :=
    \overline X'\times_{\overline Y}Y_0.
\]
The open immersion
\[
    \overline X\hookrightarrow W
\]
is proper, because \(\overline X\to Y_0\) is proper and \(W\to Y_0\) is
separated.  Hence it is closed.  It is also schematically dense, so
\[
    W\simeq\overline X.
\]
Thus the central square is Cartesian and the resulting datum refines every
original one.  For a finite diagram, compatibility with all arrows and
higher simplices follows from uniqueness of maps out of a dense common
refinement.  Hence the category is cofiltered.
\end{proof}

\begin{construction}[Lifting classical composition data]
\label{constr:composition-data-lifting}
Classical truncation defines
\[
    \tau_{f,g}:
    \operatorname{Comp}(f,g)
    \longrightarrow
    \operatorname{Comp}^{\square}_{\mathrm{cl}}(f_0,g_0).
\]

For a classical composition datum, define
\[
    \overline X
    :=
    X\amalg_{X_0}\overline X_0,
\]
\[
    \overline Y
    :=
    Y\amalg_{Y_0}\overline Y_0,
\]
and
\[
    \overline X'
    :=
    \overline X
    \amalg_{\overline X_0}
    \overline X'_0.
\]
The pushout universal properties induce
\[
    p_f:\overline X\longrightarrow Y,
    \qquad
    p_g:\overline Y\longrightarrow S,
    \qquad
    r:\overline X'\longrightarrow\overline Y.
\]
Their classical truncations are the corresponding proper classical maps,
so they belong to \(\mathcal P\).  The three horizontal maps are
quasi-compact open immersions.

The pullback
\[
    V
    :=
    \overline X'\times_{\overline Y}Y
\]
is an open spectral subscheme of \(\overline X'\), and
\[
    V_{\mathrm{cl}}
    \simeq
    \overline X'_0
    \times_{\overline Y_0}
    Y_0
    \simeq
    \overline X_0.
\]
The open subscheme \(\overline X\subseteq\overline X'\) has the same
classical open.  Hence
\[
    V\simeq\overline X,
\]
so the central square is Cartesian.  This construction is functorial in
strict refinements and defines
\[
    \lambda_{f,g}:
    \operatorname{Comp}^{\square}_{\mathrm{cl}}(f_0,g_0)
    \longrightarrow
    \operatorname{Comp}(f,g).
\]
\end{construction}

\begin{proposition}[The composition-data lifting adjunction]
\label{prop:composition-data-lifting-adjunction}
There is an adjunction
\[
    \lambda_{f,g}:
    \operatorname{Comp}^{\square}_{\mathrm{cl}}(f_0,g_0)
    \rightleftarrows
    \operatorname{Comp}(f,g):
    \tau_{f,g}.
\]
Its unit is an equivalence, and its counit is the natural proper refinement
obtained from the truncation immersions of the three compactifying spectral
schemes.
\end{proposition}

\begin{proof}
Apply the mapping-space universal property of the three nil-cobase pushouts.
Maps from the classical pieces into spectral schemes are identified, by the
adjunction between discrete enhancement and classical truncation, with maps
into the corresponding classical truncations.  After fixing the maps on
\(X\) and \(Y\), the resulting mapping space is exactly the discrete mapping
space of strict classical refinements.  Properness is detected on classical
truncations, and each Cartesian open-square condition is detected by the
classical open subscheme as in
Proposition~\ref{prop:compactification-lifting-adjunction}.  Thus there are
natural equivalences
\[
\begin{aligned}
    &\Map_{\operatorname{Comp}(f,g)}
    \bigl(
      \lambda_{f,g}(c_0),c
    \bigr) \\
    &\qquad\simeq
    \Map_{\operatorname{Comp}^{\square}_{\mathrm{cl}}(f_0,g_0)}
    \bigl(
      c_0,\tau_{f,g}(c)
    \bigr).
\end{aligned}
\]
This proves the adjunction.  The classical truncation of every lifted object
is the original classical object, giving the unit equivalence.  The maps from
the three nil-cobase pushouts to the original spectral compactifying schemes
give the counit; their classical truncations are identities, and their open
squares are Cartesian by classical-open detection.
\end{proof}

\begin{proposition}[Compactifications and composition]
\label{prop:compactification-composition-data}
Let
\[
    X\xrightarrow{f}Y\xrightarrow{g}S
\]
belong to \(\mathcal E\).  The \(\infty\)-category
\[
    \operatorname{Comp}(f,g)
\]
is nonempty and cofiltered.  Every finite diagram of composition data,
proper refinements, and coherence simplices admits a coherently compatible
refining cone.  Consequently,
\[
    \bigl|
      \operatorname{Comp}(f,g)
    \bigr|
    \simeq
    *.
\]
\end{proposition}

\begin{proof}
Nonemptiness follows by applying \(\lambda_{f,g}\) to an object supplied by
Proposition~\ref{prop:classical-composition-data-cofiltered}.

Let
\[
    D:K\longrightarrow\operatorname{Comp}(f,g)
\]
be a diagram indexed by a finite simplicial set.  By classical
cofilteredness, \(\tau_{f,g}D\) admits a cone
\[
    c_0\longrightarrow\tau_{f,g}D.
\]
Applying \(\lambda_{f,g}\) and then composing with the natural counit
\[
    \lambda_{f,g}\tau_{f,g}D
    \longrightarrow
    D
\]
gives a coherently compatible cone over \(D\).  Therefore
\(\operatorname{Comp}(f,g)\) is cofiltered, and its groupoid completion is
contractible.
\end{proof}

\begin{remark}[Composition represented by composition data]
\label{rem:composition-data-exceptional-formula}
For an object of \(\operatorname{Comp}(f,g)\), open--proper exchange for the
central Cartesian square gives
\[
    r_*k_\sharp
    \simeq
    j_{g\sharp}p_{f*}.
\]
Consequently,
\[
\begin{aligned}
    p_{g*}j_{g\sharp}p_{f*}j_{f\sharp}
    &\simeq
    p_{g*}r_*k_\sharp j_{f\sharp} \\
    &\simeq
    (p_gr)_*(kj_f)_\sharp.
\end{aligned}
\]
The right-hand side is the compactification formula associated to the
compactification
\[
    X\xrightarrow{kj_f}\overline X'\xrightarrow{p_gr}S
\]
of \(gf\).  This is the comparison used to construct composition of
exceptional direct images in the next chapter.
\end{remark}

\subsection{Finite compactification grids}

Let
\[
    X_n
    \xrightarrow{f_n}
    X_{n-1}
    \longrightarrow
    \cdots
    \longrightarrow
    X_1
    \xrightarrow{f_1}
    X_0
\]
be a finite composable string of morphisms in \(\mathcal E\).

\begin{definition}[Finite compactification grid]
\label{def:finite-compactification-grid}
A compactification grid for
\[
    (f_1,\ldots,f_n)
\]
consists of a coherently commutative triangular diagram of spectral schemes
\[
    \overline X_{a,b},
    \qquad
    0\leq a\leq b\leq n,
\]
together with the following data.

\begin{enumerate}
    \item For every \(0\leq a\leq n\), there is a specified equivalence
    \[
        \overline X_{a,a}
        \simeq
        X_a.
    \]

    \item For every \(a<b\), there is a quasi-compact open immersion
    \[
        u_{a,b}:
        \overline X_{a+1,b}
        \hookrightarrow
        \overline X_{a,b}.
    \]

    \item For every \(a<b\), there is a morphism in \(\mathcal P\)
    \[
        p_{a,b}:
        \overline X_{a,b}
        \longrightarrow
        \overline X_{a,b-1}.
    \]

    \item For every \(a+1<b\), the square
    \[
    \begin{tikzcd}[column sep=large,row sep=large]
        \overline X_{a+1,b}
            \arrow[r,hook,"u_{a,b}"]
            \arrow[d,"p_{a+1,b}"']
            &
        \overline X_{a,b}
            \arrow[d,"p_{a,b}"]
        \\
        \overline X_{a+1,b-1}
            \arrow[r,hook,"u_{a,b-1}"']
            &
        \overline X_{a,b-1}
    \end{tikzcd}
    \]
    is Cartesian.

    \item The induced morphism
    \[
        X_b
        =
        \overline X_{b,b}
        \longrightarrow
        \overline X_{a,a}
        =
        X_a
    \]
    is the specified composite
    \[
        f_{a+1}\cdots f_b.
    \]
\end{enumerate}

A refinement of compactification grids is a morphism of the entire
triangular diagrams which is the identity on the diagonal objects \(X_a\),
whose off-diagonal components belong to \(\mathcal P\), and whose squares
with every horizontal open immersion are Cartesian.

Write
\[
    \operatorname{Comp}^{\triangle}
    (f_1,\ldots,f_n)
\]
for the resulting \(\infty\)-category.  For \(n=1\), this is
\(\operatorname{Comp}(f_1)\).  For \(n=2\), it is canonically identified
with the category \(\operatorname{Comp}(f_2,f_1)\) of pairwise composition
data constructed above.
\end{definition}

\begin{construction}[Restriction of compactification grids]
\label{constr:restriction-finite-compactification-grids}
Let
\[
    \overline X_\bullet
    \in
    \operatorname{Comp}^{\triangle}
    (f_1,\ldots,f_n)
\]
be a compactification grid, and let
\[
    \phi:[m]\longrightarrow[n]
\]
be a monotone map.

Define the restricted triangular diagram
\[
    \phi^*\overline X_\bullet
\]
on objects by
\[
    \bigl(
      \phi^*\overline X
    \bigr)_{a,b}
    :=
    \overline X_{\phi(a),\phi(b)},
    \qquad
    0\leq a\leq b\leq m.
\]

For \(a<b\), define its horizontal morphism
\[
    \bigl(
      \phi^*u
    \bigr)_{a,b}:
    \overline X_{\phi(a+1),\phi(b)}
    \longrightarrow
    \overline X_{\phi(a),\phi(b)}
\]
to be the composite
\[
    \overline X_{\phi(a+1),\phi(b)}
    \longrightarrow
    \overline X_{\phi(a+1)-1,\phi(b)}
    \longrightarrow
    \cdots
    \longrightarrow
    \overline X_{\phi(a),\phi(b)}
\]
of the horizontal open immersions of
\(\overline X_\bullet\).  When
\[
    \phi(a)=\phi(a+1),
\]
this composite is the identity.

For \(a<b\), define its vertical morphism
\[
    \bigl(
      \phi^*p
    \bigr)_{a,b}:
    \overline X_{\phi(a),\phi(b)}
    \longrightarrow
    \overline X_{\phi(a),\phi(b-1)}
\]
to be the composite
\[
    \overline X_{\phi(a),\phi(b)}
    \longrightarrow
    \overline X_{\phi(a),\phi(b)-1}
    \longrightarrow
    \cdots
    \longrightarrow
    \overline X_{\phi(a),\phi(b-1)}
\]
of the proper vertical morphisms of
\(\overline X_\bullet\).  When
\[
    \phi(b-1)=\phi(b),
\]
this composite is the identity.

The horizontal morphisms of
\(\phi^*\overline X_\bullet\) are quasi-compact open immersions, and its
vertical morphisms belong to \(\mathcal P\), by closure of those classes
under composition.  Every elementary square of the restricted diagram is
Cartesian by pasting the elementary Cartesian squares of
\(\overline X_\bullet\); when an index is repeated, the corresponding
square is an identity square.

The diagonal object at \(a\) is
\[
    \bigl(
      \phi^*\overline X
    \bigr)_{a,a}
    =
    X_{\phi(a)},
\]
and the induced morphism
\[
    X_{\phi(b)}
    \longrightarrow
    X_{\phi(a)}
\]
is the specified composite of the morphisms in the original string between
those two indices.  Thus
\(\phi^*\overline X_\bullet\) is a compactification grid for the
composable string obtained from
\[
    X_n\longrightarrow\cdots\longrightarrow X_0
\]
by restriction along \(\phi\), with identity morphisms inserted whenever
\(\phi\) repeats an index.

This construction is functorial in refinements.  For composable monotone
maps
\[
    [\ell]\xrightarrow{\psi}[m]\xrightarrow{\phi}[n],
\]
associativity of composition gives a canonical coherent equivalence
\[
    \psi^*\phi^*\overline X_\bullet
    \simeq
    (\phi\psi)^*\overline X_\bullet.
\]
For the identity monotone map, the restriction is canonically the original
grid.  These equivalences satisfy the usual higher functoriality
coherences.
\end{construction}

\begin{definition}[Strict classical compactification grids]
\label{def:strict-classical-compactification-grid}
Put
\[
    X_{a,0}
    :=
    (X_a)_{\mathrm{cl}},
    \qquad
    f_{a,0}
    :=
    (f_a)_{\mathrm{cl}}.
\]
Let
\[
    \operatorname{Comp}^{\triangle,\square}_{\mathrm{cl}}
    (f_{1,0},\ldots,f_{n,0})
\]
be the ordinary category, regarded through its nerve as an
\(\infty\)-category, defined by the same triangular diagrams with ordinary
proper morphisms, quasi-compact open immersions, strict Cartesian squares,
and proper strict refinements.

A classical compactification grid is dense if every horizontal open
immersion
\[
    \overline X_{a+1,b,0}
    \hookrightarrow
    \overline X_{a,b,0}
\]
is schematically dense.
\end{definition}

\begin{lemma}[Relative common refinement of one grid cell]
\label{lem:relative-grid-cell-refinement}
Let \(\Lambda\) be a finite set.  For every \(\lambda\in\Lambda\), suppose
that there is a Cartesian square of ordinary schemes
\[
\begin{tikzcd}[column sep=large,row sep=large]
    U_\lambda
        \arrow[r,hook]
        \arrow[d]
        &
    B_\lambda
        \arrow[d]
    \\
    W_\lambda
        \arrow[r,hook]
        &
    V_\lambda,
\end{tikzcd}
\]
whose horizontal morphisms are schematically dense quasi-compact open
immersions and whose vertical morphisms are proper.

Suppose that there are compatible common dense refinements
\[
    U\longrightarrow U_\lambda,
    \qquad
    W\longrightarrow W_\lambda,
    \qquad
    V\longrightarrow V_\lambda,
\]
such that
\[
    U\longrightarrow W
\]
is proper,
\[
    W\hookrightarrow V
\]
is schematically dense and open, and every square
\[
\begin{tikzcd}[column sep=large,row sep=large]
    W \arrow[r,hook] \arrow[d]
        & V \arrow[d] \\
    W_\lambda \arrow[r,hook]
        & V_\lambda
\end{tikzcd}
\]
is Cartesian.

Then there exist a proper morphism
\[
    B\longrightarrow V,
\]
a schematically dense quasi-compact open immersion
\[
    U\hookrightarrow B,
\]
and proper morphisms
\[
    B\longrightarrow B_\lambda
\]
which refine all the displayed squares and for which
\[
    B\times_VW
    \simeq
    U.
\]
\end{lemma}

\begin{proof}
For every \(\lambda\), put
\[
    Q_\lambda
    :=
    B_\lambda\times_{V_\lambda}V.
\]
Then \(Q_\lambda\to V\) is proper.  The boundary refinement maps and the
Cartesian square for \(B_\lambda\) induce a morphism
\[
    U\longrightarrow Q_\lambda.
\]
Put
\[
    Q
    :=
    \prod_VQ_\lambda.
\]
The morphism \(Q\to V\) is proper, and the preceding maps give
\[
    U\longrightarrow Q.
\]

The composite
\[
    U\longrightarrow W\hookrightarrow V
\]
is separated and of finite type.  Choose a dense classical Nagata
compactification
\[
    U\hookrightarrow P\longrightarrow V.
\]
The pair of maps from \(U\) determines a locally closed immersion
\[
    U
    \longrightarrow
    P\times_VQ.
\]
Let \(B\) be its scheme-theoretic closure.  Then
\[
    B\longrightarrow V
\]
is proper and
\[
    U\hookrightarrow B
\]
is schematically dense and open.

Put
\[
    B_W
    :=
    B\times_VW.
\]
The immersion
\[
    U\hookrightarrow B_W
\]
is open.  It is also proper, because \(U\to W\) is proper and
\(B_W\to W\) is separated.  Hence it is closed.  Since \(U\) is
schematically dense in \(B_W\), it follows that
\[
    B_W
    \simeq
    U.
\]

The projections
\[
    B\longrightarrow Q_\lambda\longrightarrow B_\lambda
\]
are proper.  Moreover,
\[
\begin{aligned}
    B\times_{B_\lambda}U_\lambda
    &\simeq
    B\times_{V_\lambda}W_\lambda \\
    &\simeq
    B\times_VW \\
    &\simeq
    U.
\end{aligned}
\]
Thus the induced open squares are Cartesian, and \(B\) gives the required
common refinement of the cell.
\end{proof}

\begin{proposition}[Cofilteredness of strict classical compactification grids]
\label{prop:classical-compactification-grids-cofiltered}
The category
\[
    \operatorname{Comp}^{\triangle,\square}_{\mathrm{cl}}
    (f_{1,0},\ldots,f_{n,0})
\]
is nonempty and cofiltered.
\end{proposition}

\begin{proof}
We first construct one grid.  Proceed by induction on the interval length
\[
    b-a.
\]
For \(a=b\), put
\[
    \overline X_{a,a,0}
    :=
    X_{a,0}.
\]
Suppose that all cells of interval length less than \(b-a\) have been
constructed.  The composite
\[
    \overline X_{a+1,b,0}
    \longrightarrow
    \overline X_{a+1,b-1,0}
    \hookrightarrow
    \overline X_{a,b-1,0}
\]
is separated and of finite type.  Choose a dense classical Nagata
compactification
\[
    \overline X_{a+1,b,0}
    \hookrightarrow
    \overline X_{a,b,0}
    \longrightarrow
    \overline X_{a,b-1,0}.
\]
Let
\[
    W
    :=
    \overline X_{a,b,0}
    \times_{\overline X_{a,b-1,0}}
    \overline X_{a+1,b-1,0}.
\]
The induced open immersion
\[
    \overline X_{a+1,b,0}
    \hookrightarrow
    W
\]
is proper, because its source is proper over
\(\overline X_{a+1,b-1,0}\) and \(W\) is separated over that base.  Hence
it is closed.  It is schematically dense, so
\[
    W
    \simeq
    \overline X_{a+1,b,0}.
\]
Thus the new elementary square is Cartesian.  This proves nonemptiness.

Every grid admits a dense refinement.  Construct it by induction on interval
length, taking scheme-theoretic closures at length one and applying
Lemma~\ref{lem:relative-grid-cell-refinement} to each subsequent cell.

Now consider finitely many dense grids.  Again proceed by induction on
interval length.  Once common dense refinements have been constructed on the
three boundary terms of a cell, apply
Lemma~\ref{lem:relative-grid-cell-refinement} to obtain a common dense
refinement of the cell.  This constructs a common dense refinement of every
finite family of grids.

Finally, maps out of a dense grid are unique whenever they exist.  Let
\[
    r,s:
    \overline X_\bullet
    \rightrightarrows
    \overline Y_\bullet
\]
be refinements with dense source.  Induct on interval length.  On the
diagonal the maps agree.  At the cell \((a,b)\), the induction hypothesis
shows that they agree on the schematically dense open
\[
    \overline X_{a+1,b}
    \hookrightarrow
    \overline X_{a,b}.
\]
The target \(\overline Y_{a,b}\) is separated over
\(\overline Y_{a,b-1}\), because it is proper.  Hence the equalizer of the
two maps on \(\overline X_{a,b}\) is closed.  Schematic density forces the
equalizer to be the whole source, so the two maps agree.

Given a finite simplicial diagram of grids, choose a common dense refinement
of its vertices.  Compatibility with every edge and every higher simplex is
then forced by the preceding uniqueness.  Thus every finite diagram admits
a cone, and the category is cofiltered.
\end{proof}

\begin{construction}[Lifting finite classical compactification grids]
\label{constr:finite-grid-lifting}
Classical truncation defines a functor
\[
    \tau^\triangle:
    \operatorname{Comp}^{\triangle}(f_1,\ldots,f_n)
    \longrightarrow
    \operatorname{Comp}^{\triangle,\square}_{\mathrm{cl}}
    (f_{1,0},\ldots,f_{n,0}).
\]

Let
\[
    A_\bullet
    \in
    \operatorname{Comp}^{\triangle,\square}_{\mathrm{cl}}
    (f_{1,0},\ldots,f_{n,0}).
\]
Define a spectral grid \(\lambda^\triangle(A)_\bullet\) recursively on
interval length.  Put
\[
    \lambda^\triangle(A)_{a,a}
    :=
    X_a.
\]
For \(a<b\), put
\[
    \lambda^\triangle(A)_{a,b}
    :=
    \lambda^\triangle(A)_{a+1,b}
    \amalg_{A_{a+1,b}}
    A_{a,b}.
\]
Here
\[
    A_{a+1,b}
    \longrightarrow
    \lambda^\triangle(A)_{a+1,b}
\]
is the classical truncation immersion and
\[
    A_{a+1,b}
    \hookrightarrow
    A_{a,b}
\]
is the classical horizontal open immersion.

Proposition~\ref{prop:spectral-nil-cobase-change-open} gives a
quasi-compact open immersion
\[
    \lambda^\triangle(A)_{a+1,b}
    \hookrightarrow
    \lambda^\triangle(A)_{a,b}
\]
and a canonical equivalence
\[
    \bigl(
      \lambda^\triangle(A)_{a,b}
    \bigr)_{\mathrm{cl}}
    \simeq
    A_{a,b}.
\]
The compatible boundary maps induce
\[
    \lambda^\triangle(A)_{a,b}
    \longrightarrow
    \lambda^\triangle(A)_{a,b-1}.
\]
Its classical truncation is the proper morphism
\[
    A_{a,b}
    \longrightarrow
    A_{a,b-1},
\]
so it belongs to \(\mathcal P\).

The pullback
\[
    V
    :=
    \lambda^\triangle(A)_{a,b}
    \times_{\lambda^\triangle(A)_{a,b-1}}
    \lambda^\triangle(A)_{a+1,b-1}
\]
is an open spectral subscheme of \(\lambda^\triangle(A)_{a,b}\), and
\[
\begin{aligned}
    V_{\mathrm{cl}}
    &\simeq
    A_{a,b}
    \times_{A_{a,b-1}}
    A_{a+1,b-1} \\
    &\simeq
    A_{a+1,b}.
\end{aligned}
\]
The open subscheme \(\lambda^\triangle(A)_{a+1,b}\) has the same
classical open.  Uniqueness of open spectral subschemes with prescribed
classical open gives
\[
    V
    \simeq
    \lambda^\triangle(A)_{a+1,b}.
\]
Thus every elementary square is Cartesian.

Strict classical refinements induce spectral refinements through the pushout
universal property.  This defines a functor
\[
    \lambda^\triangle:
    \operatorname{Comp}^{\triangle,\square}_{\mathrm{cl}}
    (f_{1,0},\ldots,f_{n,0})
    \longrightarrow
    \operatorname{Comp}^{\triangle}(f_1,\ldots,f_n).
\]
\end{construction}

\begin{proposition}[The finite-grid lifting adjunction]
\label{prop:finite-grid-lifting-adjunction}
There is an adjunction
\[
    \lambda^\triangle:
    \operatorname{Comp}^{\triangle,\square}_{\mathrm{cl}}
    (f_{1,0},\ldots,f_{n,0})
    \rightleftarrows
    \operatorname{Comp}^{\triangle}(f_1,\ldots,f_n):
    \tau^\triangle.
\]
Its unit is an equivalence.  Its counit is the natural proper refinement
induced cellwise by the open boundary maps and the classical truncation
immersions.
\end{proposition}

\begin{proof}
Let \(A_\bullet\) be a strict classical grid and let \(C_\bullet\) be a
spectral grid.  We construct the mapping-space equivalence by induction on
interval length.  At a cell \((a,b)\), the pushout universal property gives
\[
\begin{aligned}
    \Map
    \bigl(
      \lambda^\triangle(A)_{a,b},
      C_{a,b}
    \bigr)
    \simeq{}
    &\Map
    \bigl(
      \lambda^\triangle(A)_{a+1,b},
      C_{a,b}
    \bigr)
    \\
    &\times_{
      \Map(A_{a+1,b},C_{a,b})
    }
    \Map(A_{a,b},C_{a,b}).
\end{aligned}
\]
The first factor is fixed by the already-constructed map on the horizontal
boundary.  Since \(A_{a,b}\) is ordinary, discrete enhancement and classical
truncation give
\[
    \Map_{\operatorname{Sch}^{\mathrm{sp}}}
    (A_{a,b},C_{a,b})
    \simeq
    \Map_{\operatorname{Sch}}
    (A_{a,b},(C_{a,b})_{\mathrm{cl}}).
\]
The compatibility condition over \(A_{a+1,b}\), together with naturality in
the vertical direction, is exactly the strict classical grid condition.
Taking the finite limit over all cells gives a natural equivalence
\[
\begin{aligned}
    &\Map_{\operatorname{Comp}^{\triangle}}
    \bigl(
      \lambda^\triangle(A),
      C
    \bigr) \\
    &\qquad\simeq
    \Map_{\operatorname{Comp}^{\triangle,\square}_{\mathrm{cl}}}
    \bigl(
      A,
      \tau^\triangle(C)
    \bigr).
\end{aligned}
\]
Properness and the Cartesian open-square conditions are detected on
classical truncations, so this equivalence restricts to the refinement
mapping spaces.  It proves the adjunction.

The equivalences
\[
    \bigl(
      \lambda^\triangle(A)_{a,b}
    \bigr)_{\mathrm{cl}}
    \simeq
    A_{a,b}
\]
identify the unit with the identity grid.  For a spectral grid \(C\), the
pushout universal properties give cellwise maps
\[
    \lambda^\triangle\tau^\triangle(C)_{a,b}
    \longrightarrow
    C_{a,b}.
\]
Their classical truncations are identities, so they belong to
\(\mathcal P\).  Their horizontal open squares are Cartesian by
classical-open detection.  Hence they assemble into the stated counit
refinement.
\end{proof}

\begin{theorem}[Contractibility of finite compactification grids]
\label{thm:finite-compactification-grids}
For every finite composable string
\[
    X_n
    \xrightarrow{f_n}
    X_{n-1}
    \longrightarrow
    \cdots
    \longrightarrow
    X_1
    \xrightarrow{f_1}
    X_0
\]
in \(\mathcal E\), the \(\infty\)-category
\[
    \operatorname{Comp}^{\triangle}(f_1,\ldots,f_n)
\]
is nonempty and cofiltered.  Consequently,
\[
    \left|
      \operatorname{Comp}^{\triangle}(f_1,\ldots,f_n)
    \right|
    \simeq
    *.
\]
Every finite diagram of compactification grids, refinements, and coherence
simplices admits a coherently compatible refining cone.
\end{theorem}

\begin{proof}
Nonemptiness follows by applying \(\lambda^\triangle\) to a strict classical
grid supplied by
Proposition~\ref{prop:classical-compactification-grids-cofiltered}.

Let
\[
    D:K
    \longrightarrow
    \operatorname{Comp}^{\triangle}(f_1,\ldots,f_n)
\]
be indexed by a finite simplicial set.  Classical cofilteredness gives a
cone
\[
    A
    \longrightarrow
    \tau^\triangle D.
\]
Applying \(\lambda^\triangle\) and then the natural counit of
Proposition~\ref{prop:finite-grid-lifting-adjunction} gives a cone
\[
    \lambda^\triangle(A)
    \longrightarrow
    \lambda^\triangle\tau^\triangle D
    \longrightarrow
    D.
\]
Thus every finite diagram admits a coherently compatible refining cone, so
the spectral grid category is cofiltered.  Its groupoid completion is
therefore contractible.
\end{proof}

\begin{proposition}[Base change of finite compactification grids]
\label{prop:finite-grid-base-change}
Let
\[
    g:T\longrightarrow X_0
\]
be arbitrary, and write
\[
    X_{a,T}
    :=
    X_a\times_{X_0}T
\]
for the base-changed composable string.  Cellwise base change defines a
functor
\[
    g^*:
    \operatorname{Comp}^{\triangle}(f_1,\ldots,f_n)
    \longrightarrow
    \operatorname{Comp}^{\triangle}(f_{1,T},\ldots,f_{n,T}).
\]
It preserves refinements and coherently compatible refining cones and is
compatible with composition in \(g\).  For every monotone map
\[
    \phi:[m]\longrightarrow[n],
\]
there is a canonical coherent equivalence
\[
    \phi^*(g^*\overline X_\bullet)
    \simeq
    g_\phi^*(\phi^*\overline X_\bullet),
\]
where
\[
    g_\phi:
    X_{\phi(0)}\times_{X_0}T
    \longrightarrow
    X_{\phi(0)}
\]
is the induced base-change morphism.
\end{proposition}

\begin{proof}
For a compactification grid \(\overline X_\bullet\), put
\[
    (g^*\overline X)_{a,b}
    :=
    \overline X_{a,b}\times_{X_0}T.
\]
Arbitrary base change preserves quasi-compact open immersions, morphisms in
\(\mathcal P\), and Cartesian squares.  It therefore carries grids to grids
and refinements to refinements.  Base change of a refining cone is again a
refining cone.  Compatibility with composition in \(g\) is the coherence of
spectral fibre products.  For a monotone map
\[
    \phi:[m]\longrightarrow[n],
\]
both
\[
    \phi^*(g^*\overline X_\bullet)
\]
and
\[
    g_\phi^*(\phi^*\overline X_\bullet)
\]
have cell \((a,b)\) canonically equivalent to
\[
    \overline X_{\phi(a),\phi(b)}
    \times_{X_0}T.
\]
The coherence of spectral fibre products gives the asserted equivalence and
its compatibility with composition of monotone maps.
\end{proof}

\begin{proposition}[Composition represented by finite grids]
\label{prop:finite-grid-composition-formula}
Let \(\overline X_\bullet\) be a compactification grid.  For \(a<b\), let
\[
    j_{a,b}:
    X_b
    \hookrightarrow
    \overline X_{a,b}
\]
be the composite horizontal open immersion and let
\[
    q_{a,b}:
    \overline X_{a,b}
    \longrightarrow
    X_a
\]
be the composite proper morphism.  Put
\[
    F_{a,b}
    :=
    q_{a,b*}j_{a,b\sharp}.
\]

For every \(a<k<b\), the inverse of open--proper exchange gives a canonical
equivalence
\[
    F_{a,k}F_{k,b}
    \xrightarrow{\simeq}
    F_{a,b}.
\]
For every subdivision
\[
    a=i_0<i_1<\cdots<i_r=b,
\]
all parenthesized composites
\[
    F_{i_0,i_1}
    F_{i_1,i_2}
    \cdots
    F_{i_{r-1},i_r}
\]
are coherently identified with \(F_{a,b}\).  These equivalences are
compatible with refinements of grids and with restriction along monotone
maps of finite ordinals.
\end{proposition}

\begin{proof}
Pasting the elementary Cartesian squares of the grid gives a Cartesian
rectangle
\[
\begin{tikzcd}[column sep=large,row sep=large]
    \overline X_{k,b}
        \arrow[r,hook,"u"]
        \arrow[d,"q_{k,b}"']
        &
    \overline X_{a,b}
        \arrow[d,"r"]
    \\
    X_k
        \arrow[r,hook,"j_{a,k}"']
        &
    \overline X_{a,k},
\end{tikzcd}
\]
where \(r\) is proper.  Applied to this square,
Theorem~\ref{thm:open-proper-exchange} gives the canonical equivalence
\[
    \operatorname{Ex}^{\sharp}_{*}:
    r_*u_\sharp
    \xrightarrow{\simeq}
    j_{a,k\sharp}q_{k,b*}.
\]
The comparison required for composition is its inverse:
\[
    \bigl(
      \operatorname{Ex}^{\sharp}_{*}
    \bigr)^{-1}:
    j_{a,k\sharp}q_{k,b*}
    \xrightarrow{\simeq}
    r_*u_\sharp.
\]

Using this inverse exchange equivalence, we obtain
\[
\begin{aligned}
    F_{a,k}F_{k,b}
    &=
    q_{a,k*}j_{a,k\sharp}
    q_{k,b*}j_{k,b\sharp} \\
    &\xrightarrow{\simeq}
    q_{a,k*}r_*u_\sharp j_{k,b\sharp} \\
    &\simeq
    q_{a,b*}j_{a,b\sharp} \\
    &=
    F_{a,b}.
\end{aligned}
\]
Here the second equivalence uses the canonical composition equivalences
\[
    q_{a,k}r
    \simeq
    q_{a,b}
\]
and
\[
    uj_{k,b}
    \simeq
    j_{a,b}.
\]

The composition compatibility of open--proper exchange identifies the
exchange transformation associated to a pasted rectangle with the composite
of the exchange transformations associated to its constituent rectangles.
Taking inverses reverses the order of this pasted composite in precisely the
way required by functor composition.  Hence the displayed identifications
satisfy all associativity and higher pasting coherences.

Compatibility with refinements follows from the composition coherence of
the proper support and open--proper exchange equivalences.  Namely, a
refinement gives a map between the corresponding Cartesian rectangles, and
naturality of open--proper exchange identifies the resulting comparison
maps.

Restriction along a monotone map is the construction of
Construction~\ref{constr:restriction-finite-compactification-grids}.  Its
horizontal and vertical arrows are composites of those in the original
grid, and its Cartesian rectangles are obtained by pasting the original
Cartesian squares.  Naturality and composition compatibility of
open--proper exchange therefore give the asserted face, degeneracy, and
higher restriction compatibilities.
\end{proof}

\section{The Voevodsky conditions}
\label{sec:voevodsky-conditions}

The preceding results may be collected as the geometric conditions required
to construct exceptional direct image by compactification.

\begin{theorem}[Voevodsky conditions for stable brave new motives]
\label{thm:voevodsky-conditions}
The stable coefficient system
\[
    \SH^*:
    \bigl(
      \operatorname{Sch}^{\mathrm{sp}}_{\mathrm{qcqs}}
    \bigr)^{\mathrm{op}}
    \longrightarrow
    \CAlg(\Pr^L_{\mathrm{st}})
\]
has the following properties.
\begin{enumerate}
    \item \textbf{Localization.}
    For every complementary closed--open pair
    \[
        Z\xrightarrow{i}S\xleftarrow{j}U,
    \]
    there are exact triangles
    \[
        j_\sharp j^*E
        \longrightarrow
        E
        \longrightarrow
        i_*i^*E
    \]
    and
    \[
        i_*i^!E
        \longrightarrow
        E
        \longrightarrow
        j_*j^*E.
    \]

    \item \textbf{Smooth exchange.}
    Smooth base change and the smooth projection formula hold for every
    smooth separated morphism of finite presentation.

    \item \textbf{Smooth proper ambidexterity.}
    For every smooth proper morphism \(p\),
    \[
        p_\sharp\Sigma^{-T_p}
        \simeq
        p_*.
    \]

    \item \textbf{Cocontinuity.}
    For every morphism \(f\), the direct image \(f_*\) preserves small
    colimits.

    \item \textbf{Proper exchange.}
    Proper base change and the proper projection formula hold for every
    proper morphism.

    \item \textbf{Open--proper exchange and support.}
    For every Cartesian square with a proper vertical morphism and a
    quasi-compact open horizontal morphism,
    \[
        p_*j'_\sharp
        \simeq
        j_\sharp q_*.
    \]
    In particular, proper direct image carries extension by zero from the
    source open to extension by zero from the target open in every proper
    excision square.

    \item \textbf{Compactification.}
    Every separated morphism of finite presentation admits a compactification,
    and its \(\infty\)-category of compactifications is nonempty and
    cofiltered.  More generally, the \(\infty\)-category of compactification
    grids for every finite composable string in \(\mathcal E\) is nonempty
    and cofiltered.
\end{enumerate}

The localization, exchange, projection, support, and ambidexterity
transformations satisfy the usual composition and base-change pasting laws.
Virtual Thom twists commute with inverse image.  If
\[
    L\in\Pic(\SH(S))
\]
is a virtual Thom object on the target, then the smooth, closed, and proper
projection formulas give canonical equivalences
\[
    p_\sharp\bigl(A\otimes_Xp^*L\bigr)
    \simeq
    p_\sharp A\otimes_SL,
\]
\[
    i_*\bigl(A\otimes_Zi^*L\bigr)
    \simeq
    i_*A\otimes_SL,
\]
and
\[
    f_*\bigl(A\otimes_Xf^*L\bigr)
    \simeq
    f_*A\otimes_SL,
\]
respectively.
For every morphism \(f\), arbitrary direct image satisfies the internal-Hom
equivalence of
Proposition~\ref{prop:arbitrary-direct-image-internal-hom}.  Finite
compactification-grid and refinement diagrams admit coherently compatible
refining cones.
\end{theorem}

\begin{proof}
Localization and smooth exchange were proved in
Chapter~\ref{chap:stable-brave-new-motivic}.  Smooth proper ambidexterity is
Theorem~\ref{thm:smooth-proper-ambidexterity}.  Proper base change and the
proper projection formula are
Theorems~\ref{thm:proper-base-change} and
\ref{thm:proper-projection-formula}.  Cocontinuity is
Theorem~\ref{thm:arbitrary-direct-image-cocontinuous}.  Open--proper exchange
is Theorem~\ref{thm:open-proper-exchange}, and proper support is
Corollary~\ref{thm:proper-support}.  Compactification is
Theorem~\ref{thm:spectral-nagata-compactification}, while cofilteredness and
contractibility are
Theorem~\ref{thm:compactification-category-contractible}.

The composition and base-change compatibilities follow from the coherent
construction of the canonical Beck--Chevalley, projection, purity, and
ambidexterity transformations and from the calculus of adjoint mates.
Compatibility with virtual Thom twists follows from pullback naturality of
the motivic \(J\)-homomorphism and the appropriate projection formulas.
The internal-Hom statement is
Proposition~\ref{prop:arbitrary-direct-image-internal-hom}.  Coherent
refinements for single compactifications are supplied by
Proposition~\ref{prop:finite-compactification-diagram-refinement}, and the
finite composition coherences are supplied by
Theorem~\ref{thm:finite-compactification-grids},
Proposition~\ref{prop:finite-grid-base-change}, and
Proposition~\ref{prop:finite-grid-composition-formula}.
\end{proof}

\section{Classical comparison of the proper package}
\label{sec:proper-classical-comparison}

\begin{theorem}[Comparison of projective and proper geometry]
\label{thm:proper-package-comparison}
Let
\[
    f:X\longrightarrow S
\]
be a morphism of quasi-compact quasi-separated spectral schemes.  Stable
classical comparison identifies the structures of this chapter with their
ordinary counterparts as follows.

\begin{enumerate}
    \item For every \(f\), there is a canonical equivalence
    \[
        \Theta_S^{\mathrm{st}}f_*
        \simeq
        (f_{\mathrm{cl}})_*
        \Theta_X^{\mathrm{st}}.
    \]

    \item If \(f\) is smooth and proper, comparison carries
    \[
        f_\sharp\Sigma^{-T_f}
        \xrightarrow{\simeq}
        f_*
    \]
    to ordinary smooth proper ambidexterity.

    \item If \(f\) is proper, comparison carries spectral proper base change,
    the proper projection formula, open--proper exchange, and the proper
    support equivalence to their ordinary counterparts.

    \item Classical truncation carries every spectral compactification
    \[
        X\hookrightarrow\overline X\longrightarrow S
    \]
    to a classical compactification
    \[
        X_{\mathrm{cl}}
        \hookrightarrow
        \overline X_{\mathrm{cl}}
        \longrightarrow
        S_{\mathrm{cl}}.
    \]
    It also carries proper refinements and finite compactification grids to
    their classical counterparts and is compatible with base change and
    finite refinement diagrams.
\end{enumerate}
\end{theorem}

\begin{proof}
The first assertion is the arbitrary-direct-image compatibility of stable
classical comparison established in
Chapter~\ref{chap:stable-brave-new-motivic}.  It is obtained by taking right
adjoints of the comparison equivalence for inverse image and does not require
\(f\) to be proper.

Suppose that \(f\) is smooth and proper.  The construction of the
ambidexterity transformation uses smooth base change, the diagonal
adjunction, the coefficientwise relative closed-purity transformation and
its diagonal Thom mate, the smooth and closed projection formulas, and the
tangent Thom twist.  Stable comparison preserves each of
these ingredients and their adjunction data.  It therefore carries
\[
    f_\sharp\Sigma^{-T_f}
    \longrightarrow
    f_*
\]
to the ordinary smooth-proper ambidexterity transformation.

Now suppose that \(f\) is proper.  The proper Beck--Chevalley and projection
transformations are characterized as adjoint mates of, respectively, the
inverse-image composition equivalence and the counit of
\[
    f^*\dashv f_*.
\]
Stable comparison preserves the adjunction, its unit and counit, Cartesian
composition, and the symmetric monoidal structures.  It therefore carries
the spectral transformations to the ordinary proper base-change and
projection transformations.

Open--proper exchange is constructed from proper base change,
open--closed orthogonality, and stable recollement.  Stable comparison
preserves each of these constructions and hence carries the spectral exchange
map, and its proper-support specialization, to their ordinary counterparts.

Finally, classical truncation preserves quasi-compact open immersions,
carries morphisms in \(\mathcal P\) to proper morphisms, and preserves fibre
products and composition.  It therefore carries spectral compactifications
to classical compactifications and carries proper refinements, finite
compactification grids, base changes, and finite refinement diagrams to
their classical counterparts.
\end{proof}

\section{The proper geometric package}
\label{sec:proper-geometric-package}

\begin{theorem}[Proper morphisms, ambidexterity, and compactification]
\label{thm:proper-geometric-package}
For every quasi-compact quasi-separated spectral scheme \(S\), the stable
brave new motivic category \(\SH(S)\) has the following additional structures
and properties.
\begin{enumerate}
    \item Every smooth proper morphism
    \[
        p:X\longrightarrow S
    \]
    satisfies canonical ambidexterity
    \[
        p_\sharp\Sigma^{-T_p}
        \simeq
        p_*.
    \]
    In particular,
    \[
        \bigl(p_\sharp\mathbf1_X\bigr)^\vee
        \simeq
        p_\sharp\Sigma^{-T_p}\mathbf1_X.
    \]

    \item For every Cartesian square with \(f\) proper,
    \[
    \begin{tikzcd}[column sep=large,row sep=large]
        X_T \arrow[r,"g'"] \arrow[d,"f'"']
            & X \arrow[d,"f"] \\
        T \arrow[r,"g"'] & S,
    \end{tikzcd}
    \]
    the canonical transformation is an equivalence:
    \[
        g^*f_*
        \simeq
        f'_*(g')^*.
    \]

    \item For every morphism \(f\), the functor \(f_*\) preserves small
    colimits.  In particular, for every compactification
    \[
        X\xrightarrow{j}\overline X\xrightarrow{p}S,
    \]
    the functor \(p_*j_\sharp\) is colimit-preserving.

    \item Proper direct image satisfies the projection formula
    \[
        f_*(A\otimes_Xf^*B)
        \simeq
        f_*A\otimes_SB.
    \]

    \item Proper direct image satisfies open--proper exchange and hence the
    support property for every proper excision square.

    \item Every separated morphism of finite presentation
    \[
        f:X\longrightarrow S
    \]
    admits a compactification
    \[
        X\xrightarrow{j}\overline X\xrightarrow{p}S.
    \]
    The \(\infty\)-category of all such compactifications is nonempty and
    cofiltered, and hence has contractible groupoid completion.  Base change
    preserves compactifications and refinements.  For every finite composable
    string in \(\mathcal E\), the \(\infty\)-category of compactification
    grids is nonempty and cofiltered, and every finite diagram of grids,
    refinements, and coherence simplices admits a coherently compatible
    refining cone.
\end{enumerate}
\end{theorem}

\begin{proof}
This is the combination of
Theorems~\ref{thm:motivic-atiyah-duality},
\ref{thm:smooth-proper-ambidexterity},
\ref{thm:proper-base-change},
\ref{thm:arbitrary-direct-image-cocontinuous},
\ref{thm:proper-projection-formula},
\ref{thm:open-proper-exchange},
\ref{thm:proper-support},
\ref{thm:spectral-nagata-compactification}, and
\ref{thm:compactification-category-contractible}, together with
Proposition~\ref{prop:finite-compactification-diagram-refinement},
Proposition~\ref{prop:compactification-composition-data},
Theorem~\ref{thm:finite-compactification-grids}, and
Proposition~\ref{prop:finite-grid-base-change}.
\end{proof}

The chapter has supplied all proper geometric input required for exceptional
functoriality, but it has not yet selected a general exceptional direct image
or constructed its right adjoint.  In particular, for a general separated
morphism of finite presentation
\[
    f:X\longrightarrow S,
\]
no functor denoted \(f_!\) or \(f^!\) has yet been fixed.

For every compactification
\[
    c=
    \bigl(
      X\xrightarrow{j}\overline X\xrightarrow{p}S
    \bigr),
\]
Proposition~\ref{prop:compactification-direct-image-diagram}
associates the functor
\[
    \Phi_f(c)
    =
    p_*j_\sharp.
\]
Open--proper exchange identifies these functors under every proper
refinement, and the resulting comparison equivalences satisfy all higher
composition coherences.  Each functor \(p_*j_\sharp\) preserves small
colimits.  Since
\[
    \bigl|
      \operatorname{Comp}(f)
    \bigr|
    \simeq
    *,
\]
the colimit
\[
    \operatorname*{colim}_{c\in\operatorname{Comp}(f)}
    \Phi_f(c)
\]
is canonically equivalent to \(\Phi_f(c)\) for every compactification
\(c\), through a contractible space of compatible equivalences.

The next chapter denotes this colimit functor by
\[
    f_!.
\]
It proves functoriality under composition and base change using the
contractible compactification-grid categories constructed above.  Their
restriction maps along monotone maps of finite ordinals supply identity,
associativity, and all higher composition coherences.  Since \(f_!\) is
colimit-preserving, its right adjoint defines
\[
    f^!,
\]
completing the six-operation formalism.
